% Modernised reading — fonds Alexandre Grothendieck, shelfmark 156-6, pages 1-93.
%
% This is an INTERPRETATION, not a transcription. Everything the critical
% apparatus carried moves into footnotes. In any dispute of substance the
% transcription governs: whoever cites Grothendieck must cite the
% batch-NN.fr.tex files of this folder.
%
% Pass: Opus 5.5 (claude-opus-5-5), 2026-10-10 — first pass, unchecked against
% the pages by a human. Made on Opus 5.5, not on the Opus 5 against which the
% edition's standard was calibrated (folder 115); the two were never compared.
% The folder was read whole: five batches, pages 1-93, all mathematics, his
% pagination coinciding with the archivists'. Dates in his hand: 18 June 1986
% (p. 1), 19 June (p. 20), 20 June (additions pp. 10, 13; p. 34), 21 June
% (p. 64), 22 June (p. 87) — two days past the catalogue's 18-20/06/1986.
%
% What the folder is. Chapter VI (« GF VI ») of Vers une géométrie des formes,
% « Analysis situs (deuxième mouture) »: the second draft of chapter IV
% (folder 156-4, 10-14 June), using the set-theoretic figures of chapter V
% (folder 156-5, 14-15 June). It axiomatises an « atelier » (𝔉, ≤, ≪): figures,
% sub-figure order, refinement order, with multistrates 𝓜 (irreducibles) and
% lieux ℒ (≪-minimal non-empty figures); At 1-At 14, their description inside
% 𝓜 (At*), the ensemblist case (Atens), sub-ateliers (Sousat), and defines the
% subdivision ≼ from ≪ and disjointness (p. 85), where GF IV took it as
% primitive.
%
% Notation fixed for the whole document. 𝔉, 𝓜, ℒ; F̃ = {X ∈ 𝓜 | X ≤ F}; ≤, ≪,
% ≼, ◁ (incidence), ∥ (disjointness), ≬ (\between, compatibility), ∦
% (\nparallel, his barred ∥); Omb/Multomb (large), omb/multomb (small);
% interior refinement ≪° (\overset{\circ}{\ll}); ∂X; F'|G the induced
% refinement (his F'.G, p. 66); Fig(L) the set-theoretic figures of chapter V;
% 𝔓_f closed (= down-closed) parts. « Figure » alone is an element of 𝔉;
% « figure ensembliste » an element of a Fig(L); likewise « strate » of F (an
% element of F̃) vs member of a figure ensembliste. His axiom numbers are
% kept; the « C 4, C 5, C 6 » cited on pp. 19-27 are cited by the At numbers
% they became.
%
% Substantive departures from the pages, each footnoted where it occurs:
% — p. 3 « Dépl(F) » for Fig_𝓜(F); p. 4 (1.7)-(1.8) written as the analogues
%   of (1.4)-(1.5).
% — p. 11 ◁ is transitive, not an order; ⊴ is; « F = F » read F = G.
% — p. 11 (1.23): F̃ ⊂ G̃ ⇒ F ≤ G needs At 2; supplied.
% — p. 12 (1.24 b), marked « Vérifier »: only ⇒ is asserted here, and only
%   once (1.58) is available.
% — p. 12 (1.27): ⇐ is the axiom At 6 (p. 21), said.
% — p. 18 (1.35)-(1.36): the bijection needs every down-set of every 𝓜_{≤X}
%   to be finitely generated (e.g. 𝓜_{≤X} finite); supplied. « relation de
%   disjonction » there read « compatibilité ».
% — p. 20: « ce qui implique sans doute pas » read as a negation; flagged.
% — p. 49-50: in a « trivial » atelier, At 8 forces 𝓜 to be discrete (the page
%   says « des plus restrictives »).
% — p. 49: the Lemma stated in its exact form (I' ordered; maps preserving
%   arbitrary unions and intersections); the sobrification remark is ours.
% — p. 52: of the equivalence « plongement ⟺ a'), b') » the page proves the
%   reflection of ≤ and injectivity only; said.
% — p. 56: the Scholie needs « multistrates pairwise ≬_𝓜 »; supplied.
% — p. 58 (1.57): « X ∩ Y empty or a multistrate »; F̃ ∪ {∅} has binary infs.
% — p. 62: (1.58) from At 9 alone needs the lieux of disjoint multistrates to
%   be pairwise disjoint, not shown there; the derivation of p. 70 (At 11) is
%   complete.
% — p. 63 At 9'', p. 78 At 9_Σ: the case X ∩ Y = ∅ (resp. ∅ ∈ Σ) must be
%   admitted for « At 9'' ⇒ At 9 » to hold; supplied.
% — p. 64 At 10: stated in a slanted note read in fragments; reconstructed
%   from its uses (pp. 72, 83) and from Cont 7 of GF IV; the reconstruction is
%   ours.
% — p. 90 Atens 4: conclusion X' ≬ Y' (page: X ≬ Y); p. 93: At 10 passes
%   from 𝔉 to 𝔉' (page: the other way).
%
% Modern names given and footnoted as ours: éléments complètement
% sup-premiers / sup-irréductibles, représentation de Birkhoff (1937),
% topologie d'Alexandrov (1937), parties commençantes, complexe simplicial
% abstrait, correspondance de Raney (1952) and espaces sobres (SGA 4),
% condition de faisceau / limite sur les strates, complexes cellulaires
% réguliers.
%
% What the folder announces and never establishes: (1.24 b); the completeness
% of the At* list of the Théorème-Scholie (p. 33, « J'espère que je n'ai pas
% oublié d'axiome ! »); the commutativity of (1.42) except through At 7; the
% uniqueness of the canonical subdivision of At 14 bis; whether (i)-(iii) of
% p. 83 imply (iv); the « saturation » of an atelier (pp. 80, 88), announced
% and not begun.
\documentclass[11pt,a4paper]{article}
\input{../preamble/grothendieck}

\folder{156-6}
\pages{1}{93}
\dating{1986}
\watermark{Édition de démonstration\\interprétation personnelle de l'œuvre}
\foldertitle{[Chapitre] VI. Analysis situs (deuxième mouture) : notes manuscrites (18-20/06/1986) — lecture modernisée du dossier entier}

\begin{document}

\section*{Résumé}

\begin{resume}
Un polyèdre découpé en triangles, une carte découpée en départements puis en
communes, une boule dont on a marqué le bord : ce sont des \emph{figures},
c'est-à-dire des formes munies d'un découpage en morceaux. Ces
quatre-vingt-treize feuillets, écrits du 18 au 22 juin 1986, cherchent les
règles du jeu qui permettent de parler de telles figures sans partir des
points d'un espace : ce que Grothendieck appelle un \emph{atelier}, ou une
« algèbre de figures ».

Le problème est celui que l'\emph{Esquisse d'un programme} avait posé deux ans
plus tôt : la topologie usuelle autorise des espaces trop sauvages pour la
géométrie, et l'on voudrait une topologie « modérée » dont les objets de base
soient des morceaux bien découpés. Mais au lieu de restreindre les espaces, il
prend pour donnée primitive l'ensemble des figures et deux manières de les
comparer. Une figure peut être une \emph{sous-figure} d'une autre --- on garde
certains morceaux, comme on extrait quelques départements d'une carte de
France. Elle peut aussi en être un \emph{raffinement} --- elle redécoupe plus
finement une partie du territoire, comme une carte des communes d'une région.
Les morceaux indécomposables sont les \emph{multistrates} ; les plus petites
figures possibles, celles qu'aucun raffinement ne rend plus fines, sont les
\emph{lieux}, qui jouent le rôle des points.

L'idée qui organise le dossier est celle d'\emph{ombre}. L'ombre d'une figure
est l'ensemble des lieux qu'elle recouvre ; sa \emph{multiombre} est la
famille des ombres de ses morceaux. Une carte ne se reconstitue pas à partir de
la seule région qu'elle couvre, mais elle se reconstitue souvent à partir du
dessin de chacun de ses départements : c'est ce que les feuillets appellent un
atelier \emph{ensembliste}. Viennent ensuite des axiomes de plus en plus
géométriques : qu'il y ait assez de lieux, que deux figures soient disjointes
quand leurs lieux le sont, que des raffinements donnés morceau par morceau et
d'accord sur les recouvrements se recollent en un raffinement de l'ensemble,
comme on recolle des cartes locales. Les derniers feuillets définissent enfin la
\emph{subdivision} : un raffinement qui redécoupe tout le territoire sans en
rien perdre.

Le dossier est une \emph{deuxième mouture}, recommencée huit jours après la
première, et l'on y voit le travail se faire. Il compare trois choix de notions
primitives et en retient un parce que « techniquement » c'est le plus simple,
en notant que les axiomes les plus simples ne sont pas les plus transparents.
Le 19 juin il change d'avis sur le sens même du raffinement ; le 20 il écrit
en marge d'une conclusion des jours précédents « Non, ça ne suffit pas » ; il
clôt une liste d'axiomes par « J'espère que je n'ai pas oublié d'axiome ! ».
Les noms à chercher ensuite sont ceux de la théorie des ensembles ordonnés
--- éléments irréductibles, parties commençantes, représentation de
Birkhoff, topologie d'Alexandrov ---, des complexes simpliciaux et
cellulaires, et des faisceaux.

\keywords{join-irreducible element, Birkhoff representation, Alexandrov topology, down-set, sober space, abstract simplicial complex, regular cell complex, subdivision, refinement, sheaf condition, tame topology}
\end{resume}

\section*{Le fil du dossier, et les conventions}

\begin{enumerate}
\item[1.] \emph{L'atelier : lieux, multistrates, figures, ombres} (pages 1 à 6),
le 18 juin.
\item[2.] \emph{Six relations, et lesquelles suffisent} (pages 6 à 14).
\item[3.] \emph{L'atelier vu dans les multistrates} (pages 14 à 18).
\item[4.] \emph{Le raffinement par les grandes ombres} (pages 18 à 28), avec le
changement de point de vue du 19 juin (page 20) et les axiomes At 4 à At 7.
\item[5.] \emph{Le théorème-scholie} (pages 28 à 33) : un atelier est une famille
de figures ensemblistes dans $\mathcal{M}$.
\item[6.] \emph{Petites multiombres, exemples, sous-ateliers} (pages 34 à 40),
le 20 juin.
\item[7.] \emph{Sous-ateliers spécieux et axiome des lieux} (pages 41 à 46).
\item[8.] \emph{Quasi-morphismes et plongements} (pages 47 à 52).
\item[9.] \emph{Ateliers ensemblistes} (pages 52 à 58).
\item[10.] \emph{Axiomes de disjonction} (pages 59 à 64).
\item[11.] \emph{Restreindre et recoller les raffinements} (pages 64 à 72), le
21 juin.
\item[12.] \emph{Compatibilité des raffinements, $\Sigma$-ateliers} (pages 73 à
80).
\item[13.] \emph{Subdivisions} (pages 81 à 88), le 22 juin à la page 87.
\item[14.] \emph{Le cas ensembliste, et les sous-ateliers} (pages 88 à 93).
\end{enumerate}

La page 1 porte en angle « GF VI » et le titre « Analysis situs (deuxième
mouture) ». C'est le chapitre VI d'un texte dont les chapitres I à III
(dossiers 156-1 à 156-3) ouvrent le programme et traitent notamment des failles,
des mailles et des réseaux, le
chapitre IV (dossier 156-4, du 10 au 14 juin) est la première mouture du même
« Analysis situs », et le chapitre V (dossier 156-5, « Algèbre des figures »,
14 et 15 juin) la théorie des figures ensemblistes dont celui-ci se sert
constamment. La troisième mouture suit dès le 23 juin (dossier 156-7). Ce
dossier-ci renvoie explicitement à la première mouture aux pages 13, 16, 30
et 38 ; les numéros de page qu'il y cite sont ceux de la première mouture. \emph{Deux moutures ne se
confondent pas} : là où il s'en sépare, on le signale, sans fondre les deux.

\emph{Figures ensemblistes.} Pour un ensemble $L$, $\mathrm{Fig}(L)$ est
l'ensemble des figures ensemblistes dans $L$ au sens du chapitre V : les
familles $\Phi$ de parties de $L$ telles que tout point de la réunion $|\Phi|$
soit dans un plus petit membre, que l'intersection de deux membres soit une
réunion de membres, et que chaque membre $A$ ait un intérieur $A^{\circ} = A
\smallsetminus \partial A$ non vide, $\partial A$ étant la réunion des membres
strictement contenus dans $A$ (si l'on omet cette dernière condition, la
famille est seulement \emph{préadmissible}). Les membres sont les \emph{strates}
de $\Phi$, les $A^{\circ}$ ses strates ouvertes, qui partitionnent $|\Phi|$.
$\Psi \leq \Phi$ si $\Psi$ est une sous-famille fermée de $\Phi$ ; $\Psi$ et
$\Phi$ sont compatibles si $\Psi \cup \Phi$ est une figure dans laquelle toutes
deux sont fermées ; $\Psi \ll \Phi$ si $|\Psi| \subset |\Phi|$ et si chaque
strate ouverte de $\Psi$ est contenue dans une strate ouverte de $\Phi$,
l'application $\Psi \to \Phi$ ainsi définie étant croissante\footnote{C'est la
forme que prend la définition aux pages 25 à 27 de ce dossier ; le chapitre V la
formule autrement (dossier 156-5, pages 19 et 20), par le plus petit membre de
$\Phi$ contenant un membre de $\Psi$. Les deux formes coïncident dans les cas
qu'on rencontre ici ; on ne s'en sert que dans ce sens.}.

\emph{Conventions.} $\mathfrak{F}$ est l'ensemble des figures d'un atelier,
$\mathcal{M} \subset \mathfrak{F}$ celui des multistrates, $\mathcal{L} \subset
\mathcal{M}$ celui des lieux ; une \emph{figure} sans épithète est un élément
de $\mathfrak{F}$, à ne pas confondre avec une figure ensembliste, élément d'un
$\mathrm{Fig}(L)$, et une \emph{strate} de $F \in \mathfrak{F}$ est un élément
de son déploiement $\widetilde{F} = \lbrace X \in \mathcal{M} \mid X \leq F
\rbrace$, à ne pas confondre avec un membre d'une figure ensembliste. Une partie
$A$ d'un ensemble ordonné est \emph{fermée} si $X \in A$ et $Y \leq X$
entraînent $Y \in A$ ; $\mathfrak{P}_{\mathrm{f}}$ désigne l'ensemble des
parties fermées, $\mathcal{M}_{\leq X}$ ou $\widetilde{X}$ la partie fermée
engendrée par $X$. Les relations sont $\leq$ (sous-figure), $\ll$
(raffinement), $\preccurlyeq$ (subdivision), $\triangleleft$ (incidence),
$\parallel$ (disjonction), $\between$ (compatibilité), et $\nparallel$ pour
« non disjoint ». $\mathrm{Omb}$, $\mathrm{Multomb}$ sont les grandes ombre et
multiombre, $\mathrm{omb}$, $\mathrm{multomb}$ les petites ; $X
\mathrel{\overset{\circ}{\ll}} Y$ est le raffinement intérieur ; $F'|G$ le
raffinement de $G$ induit par $F'$. Les numéros d'axiomes (At~1 à At~14, At$^{*}$,
Atens, Sousat) et de formules (1.1) à (1.61) sont les siens.

\emph{Ce que le dossier annonce et n'établit pas} : la réciproque de (1.24 b) ;
que la liste des conditions du théorème-scholie (page 33) soit complète ;
l'unicité de la subdivision canonique de At~14 bis ; si les conditions (i) à
(iii) de la page 83 entraînent (iv) ; la \emph{saturation} d'un atelier,
annoncée aux pages 80 et 88 et jamais commencée.

\pagerange{1}{6}
\section*{L'atelier : lieux, multistrates, figures, ombres (pages 1 à 6)}

Le premier alinéa, daté du 18 juin, dit l'intention. Avant de décrire une
\emph{contrée} --- le nom de la première mouture ---, il décrit une « contrée
avec notion de multistrates », la famille des multistrates jouant un peu le rôle
d'une famille d'ouverts engendrant une topologie, ou d'une famille génératrice
d'objets d'un topos. C'est l'\emph{algèbre de figures}, ou \emph{atelier}.

Un atelier comporte trois ensembles, les \emph{lieux}, les \emph{multistrates} et
les \emph{figures}, emboîtés par des injections canoniques (1.1)--(1.2) :
\[
\mathcal{L} \subset \mathcal{M} \subset \mathfrak{F} .
\]
À chaque figure $F$ est associé son \emph{déploiement} $\widetilde{F} \subset
\mathcal{M}$, ensemble de ses \emph{strates} ou multistrates \emph{incidentes}
(1.3 a). L'ensemble $\mathcal{M}$ sera ordonné par $\leq$, $\widetilde{F}$ en sera
une partie fermée, et pour une multistrate $X$ on aura (1.5)
\[
\widetilde{X} = \mathcal{M}_{\leq X} = \lbrace Y \in \mathcal{M} \mid Y \leq X
\rbrace .
\]
D'où une seconde application (1.3 b), (1.4) :
\[
F \longmapsto \mathrm{Fig}_{\mathcal{M}}(F) = \lbrace \widetilde{X} \mid X \in
\widetilde{F} \rbrace \in \mathrm{Fig}(\mathcal{M}) ,
\]
figure ensembliste dans $\mathcal{M}$ indexée fidèlement par
$\widetilde{F}$\footnote{Page 3, le second crochet écrit « via $F \mapsto
\mathrm{Dépl}(F)$ » ; d'après (1.3), $\mathrm{Dépl}(F) = \widetilde{F}$, et c'est
$\mathrm{Fig}_{\mathcal{M}}(F)$ qui est visé.}. La page 3 avertit : la figure
associée à une multistrate $X$, vue comme partie de $\mathcal{M}$, n'est pas
$\lbrace X \rbrace$ mais $\widetilde{X}$. Et la page 15 jugera cette seconde
application « un peu bidon », parce que le raffinement ne s'y lit pas.

\subsection*{Grandes et petites ombres}

L'\emph{ombre}, ou \emph{grande ombre}, d'une figure est l'ensemble des
multistrates qui la raffinent :
\[
\mathrm{Omb}(F) = \lbrace X \in \mathcal{M} \mid X \ll F \rbrace .
\]
Rien ne dit que $F \mapsto \mathrm{Omb}(F)$ soit injective. La \emph{multiombre}
est la famille des ombres des strates (1.7), (1.8) :
\[
\mathrm{Multomb}(F) = \lbrace \mathrm{Omb}(X) \mid X \in \widetilde{F} \rbrace,
\qquad \mathrm{Omb}(X) = \lbrace Y \in \mathcal{M} \mid Y \ll X \rbrace
,
\]
analogue à (1.4)--(1.5) avec $\ll$ à la place de $\leq$\footnote{La page 4 écrit
(1.7) « $\mathrm{Multomb}(F) = \lbrace X \in \widetilde{F} \mid \mathrm{Omb}\,X
\rbrace$ » et, en (1.8), « $\mathrm{Multomb}(X)$ » au membre de gauche, là où le
texte annonce $\mathrm{Omb}(X)$. On rétablit les deux formules comme les
analogues de (1.4) et (1.5), ce que dit la phrase qui les entoure et ce que la
page 19 écrit en toutes lettres.}. Elle sera une figure ensembliste dans
$\mathcal{M}$ qui détermine $F$ (pages 19 et 20).

En tirant ces familles en arrière par l'inclusion $\mathcal{L} \hookrightarrow
\mathcal{M}$ --- on intersecte chaque membre avec $\mathcal{L}$ --- on obtient
(1.9) la \emph{petite ombre} et la \emph{petite multiombre}
\[
\mathrm{omb}(F) = \mathrm{Omb}(F) \cap \mathcal{L} , \qquad \mathrm{multomb}(F) =
\lbrace \mathrm{omb}(X) \mid X \in \widetilde{F} \rbrace \in
\mathrm{Fig}(\mathcal{L}) ,
\]
la première sans raison d'être injective, la seconde y « tendant ». Les éléments
de $\mathrm{omb}(F)$ sont les « lieux de la figure ». Une figure ne se
reconstitue presque jamais à partir de son ombre ; elle se reconstitue, dans les
« ateliers ensemblistes », à partir de sa multiombre : c'est le programme des
pages 34 à 58\footnote{Que $\mathrm{multomb}(F)$ soit indexée par $\widetilde{F}$
tout entier, c'est-à-dire qu'aucune $\mathrm{omb}(X)$ ne soit vide ou
confondue avec une autre, demande l'axiome des lieux At~8, posé seulement à la
page 45.}.

\pagerange{6}{14}
\section*{Six relations, et lesquelles suffisent (pages 6 à 14)}

Sur $\mathfrak{F}$ vivent trois ordres (1.10) : $G \leq F$, $G$ est une
\emph{sous-figure} de $F$ ; $G \ll F$, $G$ est un \emph{raffinement} de $F$ ;
$G \preccurlyeq F$, $G$ est une \emph{subdivision} de $F$. Ils induisent $\leq$ et
$\ll$ sur $\mathcal{M}$ (1.11), et des ordres discrets sur $\mathcal{L}$. Une
sous-figure est un raffinement : $G \leq F \Rightarrow G \ll F$.

S'y ajoutent deux relations symétriques (1.12) : $F$ et $G$ sont
\emph{compatibles}, $F \between G$, si $\lbrace F, G \rbrace$ est majoré pour
$\leq$ ; \emph{disjoints}, $F \parallel G$, s'ils sont compatibles et si leur
borne inférieure est la figure vide $\emptyset_{\mathfrak{F}}$, plus petit élément
de $\mathfrak{F}$ (1.14), (1.15)\footnote{Il propose pour la compatibilité deux
signes, un losange partagé et le sablier rendu ici $\between$ ; il retient le
second, « plus rapide à écrire » quoique « moins suggestif ».}. La compatibilité
est réflexive, la disjonction ne l'est que pour la figure vide (1.19). La page 8
annonce, « par les axiomes », que deux figures sont disjointes si et seulement
si leurs ombres n'ont pas de lieu commun et que tout lieu de l'une est disjoint
de tout lieu de l'autre (1.13)\footnote{C'est une annonce. L'énoncé ne
deviendra exact qu'avec les axiomes de disjonction des pages 59 à 64 : le sens
« lieux disjoints $\Rightarrow$ figures disjointes » est l'axiome At~9, et le
corollaire de la page 62 donne l'équivalence dans le cas ensembliste.}. Dans un
atelier ensembliste, deux lieux sont compatibles, et disjoints si et seulement
s'ils sont distincts ; c'est ce que la page 61 posera comme Atens~2. Sur
$\mathcal{M}$, en revanche, la compatibilité est une donnée importante, qui
s'ajoute aux deux ordres : « moralement », dit la page 9, ces trois données
suffisent à retrouver l'atelier --- ce que les pages 18 et 21 préciseront.

La sixième relation est l'\emph{incidence} (1.17) : $F \mathrel{\triangleleft}
G$ si $F \in \mathcal{M}$ et $F \leq G$. Il range les six dans un diagramme
d'implications (1.18) :

\begin{tikzcd}[column sep=small]
 & & G \mathrel{\triangleleft} F \arrow[d, Rightarrow] & & \\
G \preccurlyeq F \arrow[dr, Rightarrow] & & G \leq F \arrow[dl, Rightarrow] \arrow[dr, Rightarrow] & & G \parallel F \arrow[dl, Rightarrow] \\
 & G \ll F & & G \between F &
\end{tikzcd}

La page compte quatre ordres, $\triangleleft$, $\leq$, $\preccurlyeq$, $\ll$ ; la
marge corrige : $\triangleleft$ est transitive et antisymétrique mais non
réflexive, puisque $F \mathrel{\triangleleft} F$ équivaut à $F \in \mathcal{M}$,
et l'ordre associé est $F \trianglelefteq G \iff (F \mathrel{\triangleleft} G$ ou
$F = G)$\footnote{La marge écrit la définition avec « ou $F = F$ » ; on lit
$F = G$.}.

\subsection*{Qui détermine quoi}

\begin{enumerate}
\item[a)] $\leq$ et $\triangleleft$ se déterminent mutuellement. De $\leq$ on tire
$\mathcal{M}$, l'ensemble des éléments \emph{strictement irréductibles} --- $X
\leq \sup_{i} F_{i}$ entraîne $X \leq F_{i}$ pour un $i$ (1.20) --- puis
$\triangleleft$. De $\triangleleft$ on tire $\mathcal{M} = \lbrace X \mid X
\mathrel{\triangleleft} X \rbrace$ (1.21), les déploiements (1.22), et $F \leq G
\iff \widetilde{F} \subset \widetilde{G}$ (1.23)\footnote{Le sens $\Leftarrow$ de
(1.23) suppose que toute figure soit la borne supérieure de ses strates : c'est
l'axiome At~2 de la page 16, rappelé page 41. On le suppose ici. La page 11
écrit d'ailleurs en (1.22) « $X \mathrel{\triangleleft} \mathfrak{F}$ » pour $X
\mathrel{\triangleleft} F$.}.
\item[b)] Chacun des ordres $\leq$, $\ll$ détermine la disjonction. Pour $\leq$
c'est la définition (1.24 a). Pour $\ll$ la page propose
\[
F \parallel G \iff \mathrm{Sup}_{\ll}(F,G) \text{ existe et }
\mathrm{Inf}_{\ll}(F,G) = \emptyset_{\mathfrak{F}} ,
\]
avec en marge « Vérifier »\footnote{Le sens $\Rightarrow$ est juste une fois
acquise la stabilité de la disjonction par raffinement (1.58), que les pages
62 et 70 tirent des axiomes ultérieurs : $F \vee G$ est alors aussi la borne
supérieure pour $\ll$ (par At~4), et un raffinement commun $K$ de $F$ et $G$
vérifie $K \parallel K$, donc $K = \emptyset_{\mathfrak{F}}$. La réciproque
n'est démontrée nulle part, et l'on ne l'affirme pas.} (1.24 b).
\item[c)] $\ll$ détermine $\preccurlyeq$ (1.25) :
\[
F \preccurlyeq G \iff F \ll G \text{ et, pour tout } H, \ (H \parallel F \iff H
\parallel G) .
\]
C'est, dix-huit pages avant, la définition que les pages 83 à 85 finiront par
adopter.
\item[d)] $\ll$ détermine les lieux, éléments minimaux de $\mathfrak{F}
\smallsetminus \lbrace \emptyset_{\mathfrak{F}} \rbrace$ (1.26) : un lieu est une
figure non vide dont le seul raffinement non vide est elle-même.
\item[e)] $\ll$ ne semble déterminer ni $\leq$ ni $\between$, mais $\ll$ et
$\between$ déterminent $\leq$ (1.27) :
\[
G \leq F \iff G \ll F \text{ et } G \between F .
\]
\end{enumerate}
Le sens $\Leftarrow$ de (1.27) n'est pas une conséquence de ce qui précède :
c'est l'axiome At~6 de la page 21, et son corollaire 2\footnote{Deux formules
différentes portent le numéro (1.27) sur la page 12 ; la première, biffée,
commençait une définition de $\between$ par $\ll$.}.

D'où trois présentations minimales de la structure (1.28) : par $(\leq, \ll)$,
qu'il encadre, par $(\leq, \preccurlyeq)$, par $(\ll, \between)$. La première
mouture avait choisi $(\leq, \preccurlyeq)$ et défini $\ll$ comme le préordre
qu'elles engendrent (1.29) :
\[
G \ll F \iff \text{il existe } F' \text{ avec } G \leq F' \preccurlyeq F ,
\]
« qui est bien une notion de subdivision, plutôt que de raffinement » ; mais
démontrer que c'est un ordre y était laborieux et multipliait les
axiomes\footnote{Dans la première mouture (dossier 156-4), la subdivision est
primitive et le raffinement en est déduit ; l'antisymétrie de $\ll$ y demande
l'axiome C8 (pages 82 à 85 de ce dossier-là). Le mot lu « préordre » est
douteux ; deux ordres engendrent bien un préordre.}. Il part donc maintenant de
(1.30) :
\[
\boxed{\ \leq,\ \ll \ \text{ deux ordres sur } \mathfrak{F}, \ \text{tels que}
\ F \leq G \Rightarrow F \ll G\ }
\]
la troisième présentation n'étant « qu'une curiosité ». Le 20 juin, des notes
reviennent sur ce choix : en regard de $(\leq, \preccurlyeq)$ dans (1.16), une
addition datée, en partie illisible (« Non, … nouveau … que pour $\ll$ ») ; à
côté de la même ligne de (1.28), « Non, ça ne suffit pas, d'après nouveau point
de vue… » ; et une longue note oblique, presque illisible, qui renvoie « au 20 »,
en regard de (1.29)\footnote{Ces notes sont postérieures au changement de point de vue daté du 19 juin à la
page 20, que l'on retrouve plus bas : elles disent sans doute que $(\leq,
\preccurlyeq)$ ne détermine plus $\ll$ une fois que le raffinement n'est plus
lu comme « sous-figure d'une subdivision ». C'est notre lecture.}.

La page 10 contient une remarque sur le choix de l'ensemble de base, qui vaut au
delà de ce dossier : on peut axiomatiser les ateliers sur $\mathfrak{F}$, sur
$\mathcal{M}$ ou sur $\mathcal{L}$ ; les axiomes les plus simples sont sur
$\mathfrak{F}$, mais ils y sont les moins transparents, et les plus
transparents sont sur $\mathcal{L}$, où ils sont les plus sophistiqués.

\pagerange{14}{18}
\section*{L'atelier vu dans les multistrates (pages 14 à 18)}

De $(\leq, \ll)$ on tire $\mathcal{M}$ par (1.20), puis $\mathcal{L}$ par
(1.26)\footnote{La page 14 renvoie à « (1.25) ».}, et le déploiement (1.31)
$\widetilde{F} = \lbrace X \in \mathcal{M} \mid X \leq F \rbrace$, de sorte que
$F \leq G \iff \widetilde{F} \subset \widetilde{G}$ (1.32). Le nom actuel de
$\mathcal{M}$ est celui d'ensemble des éléments \emph{complètement
sup-premiers} de l'ensemble ordonné $(\mathfrak{F}, \leq)$\footnote{La forme
(1.20), $X \leq \sup F_{i} \Rightarrow X \leq F_{i}$, est la forme « premier » ;
les pages 38 et 39 se servent de la forme « irréductible », $X = \sup F_{i}
\Rightarrow X = F_{i}$. Pour un ensemble de parties fermées stable par
réunions majorées, comme ici, les deux coïncident. La représentation d'un treillis
distributif fini par les parties commençantes de ses éléments irréductibles est
le théorème de G.~Birkhoff (1937) ; le rapprochement est le nôtre, la page ne
cite personne.}, et $F \mapsto \widetilde{F}$ identifie $\mathfrak{F}$ à un
ensemble de parties fermées de $\mathcal{M}$ (1.33), sur lequel l'ordre $\ll$
contient l'inclusion.

La donnée de $\mathfrak{F} \subset \mathfrak{P}(\mathcal{M})$ revient à celle
d'une famille \emph{préadmissible} de parties, au sens du chapitre V :
\begin{enumerate}
\item[(i)] pour tout $X \in \mathcal{M}$, il y a une plus petite $F_{X} \in
\mathfrak{F}$ contenant $X$ --- ce sera $\mathcal{M}_{\leq X}$ ;
\item[(ii)] $X \mapsto F_{X}$ est injective,
\end{enumerate}
de sorte que le préordre $X \leq Y \iff F_{X} \subset F_{Y}$ est un ordre.
Comme la partie vide est dans $\mathfrak{F}$, $\mathfrak{F}$ n'est jamais
\emph{admissible} : $\emptyset$ n'a pas d'intérieur non vide.

La première mouture avait établi en détail l'équivalence entre $(\mathfrak{F},
\leq)$ soumis à ses conditions C~1, C~2, C~3 et un ensemble ordonné
$\mathcal{M}$ muni d'un ensemble $\mathfrak{F}$ de parties fermées contenant la
partie vide et les $\mathcal{M}_{\leq X}$, stable par la condition de
compatibilité : si $\mathcal{M}_{\leq X} \cup \mathcal{M}_{\leq Y} \in
\mathfrak{F}$ pour tous $X \in F$, $Y \in G$, alors $F \cup G \in
\mathfrak{F}$\footnote{Renvoi de la page 16 à « GF IV », §§ 5, 6, 8, numéros
lus sans certitude ; la marge renomme ces conditions At~1, At~2, At~3. Leur
énoncé est rappelé page 41 : At~1, les bornes supérieures majorées existent et
$\mathfrak{F} \neq \emptyset$ ; At~2, toute figure est la borne supérieure de
ses strates ; At~3, $F \vee G$ existe dès que $X \vee Y$ existe pour toutes
strates $X$ de $F$ et $Y$ de $G$.}. « Il n'y a pas à y revenir. »

Les figures de \emph{type fini} sont celles dont le déploiement est une réunion
finie de $\mathcal{M}_{\leq X_{i}}$ (1.34). Pour elles, la chose essentielle est
la compatibilité sur $\mathcal{M}$ (1.35)--(1.36) : une relation $\between$
réflexive et symétrique, telle que
\[
X \between Y, \quad X' \leq X, \quad Y' \leq Y \ \Longrightarrow\ X' \between Y'
,
\]
les figures étant alors les réunions finies de $\mathcal{M}_{\leq X_{i}}$ avec les
$X_{i}$ deux à deux compatibles. C'est en ce sens que la page 9 disait que
$(\mathcal{M}, \leq, \between)$ équivaut « moralement » à la donnée des
figures\footnote{La correspondance biunivoque annoncée demande que toute partie
fermée d'une figure soit encore de type fini, donc que toute partie fermée de
chaque $\mathcal{M}_{\leq X}$ soit engendrée par un nombre fini d'éléments ---
par exemple que les $\mathcal{M}_{\leq X}$ soient finis. Cette hypothèse n'est
pas sur la page ; on l'ajoute. La page 18 appelle $\between$ « relation de
disjonction », alors que $\between$ note la compatibilité depuis la page 9.}.
Dans presque tous les exemples, c'est ainsi qu'on définira un atelier.

\pagerange{18}{28}
\section*{Le raffinement par les grandes ombres (pages 18 à 28)}

Reste $\ll$, qui n'a d'interprétation simple ni dans les déploiements ni dans les
$\mathrm{Fig}_{\mathcal{M}}(F)$ : c'est ici que servent les grandes ombres. La
page 19 veut que $\mathrm{Multomb}(F)$ soit une figure ensembliste dans
$\mathcal{M}$, indexée fidèlement par $\widetilde{F}$. Les conditions qu'elle
écrit à cet effet deviennent, aux pages 21 et 27, les axiomes At~4 à
At~7\footnote{Les pages 19 à 27 citent ces conditions sous des numéros C~4, C~5,
C~6 qui ne s'accordent pas avec les numéros At adoptés : « C~4 » y désigne tantôt
la condition $X \ll Y$, $X \between Y \Rightarrow X \leq Y$ (marge de la page
19), tantôt At~4 (pages 21, 22, 27). On cite partout par les numéros At, et un
C~6 biffé de la page 19 n'est pas repris.}.

\pagerange{20}{20}
\subsection*{Le changement de point de vue du 19 juin (page 20)}

Le dernier alinéa de la page 20, daté du 19 juin, abandonne la lecture du
raffinement qui guidait la première mouture, « $G$ est une sous-figure d'une
subdivision de $F$ ». Déjà dans le cas simplicial, où les multistrates ont la
combinatoire d'un simplexe, cela n'a pas l'air naturel. Pour deux simplexes
modérés $X$, $Y$ d'un $\mathbb{R}^{n}$, on dira que $Y \ll X$ si c'est le cas pour
les figures ensemblistes correspondantes --- ce qui n'entraîne sans doute pas
l'existence d'une subdivision de $X$ dont $Y$ soit une multistrate\footnote{La
page écrit « ce qui implique sans doute pas », sans le « n' » ; on lit une
négation, que demande l'abandon annoncé. La phrase s'interrompt au bas de la
page, et la page 21 commence sur une fin de phrase dont le début manque.}.

\pagerange{21}{24}
\subsection*{Les axiomes At~4 à At~6 (pages 21 à 24)}

\begin{quote}
\textbf{At~4.} Pour $F, G \in \mathfrak{F}$, $G \ll F$ si et seulement si toute
strate de $G$ raffine une strate de $F$.
\end{quote}
\begin{quote}
\textbf{At~5.} Si $X \in \mathcal{M}$ et $X \ll F$, l'ensemble $F^{X} = \lbrace Y
\in \widetilde{F} \mid X \ll Y \rbrace$, non vide par At~4, a un plus petit
élément.
\end{quote}
\begin{quote}
\textbf{At~6.} Si $X, Y \in \mathcal{M}$ sont compatibles, $X \ll Y$ entraîne $X
\leq Y$.
\end{quote}

At~4 dit que $\ll$ sur $\mathfrak{F}$ se déduit de $\ll$ sur $\mathcal{M}$ ; avec
la page 18, un atelier à figures de type fini se reconstitue donc à partir de
$(\mathcal{M}, \leq, \ll, \between)$. At~6 a deux corollaires (page 22) : sur
chaque $\widetilde{F}$, dont les éléments sont deux à deux compatibles, $\leq$ et
$\ll$ coïncident ; et pour deux figures compatibles, $F \ll G \iff F \leq G$ ---
la formule (1.27) de la page 12. Il en résulte que $\mathrm{Multomb}(F)$ est une
figure ensembliste, et même admissible : tout $X \in \mathrm{Omb}(F)$ est dans
un plus petit membre $\mathrm{Omb}(Y)$ (At~5, joint au premier corollaire), et
chaque membre a un intérieur non vide, puisque $Y \in \mathrm{Omb}(Y)^{\circ}$.

L'application $X \mapsto \mathrm{Omb}(X) = \mathcal{M}_{\ll X}$ (1.39) est
injective, $\ll$ étant un ordre : $X$ est le plus grand élément de son ombre.
$\mathrm{Multomb}(F)$ est l'image de $\widetilde{F}$ par cette injection, et
détermine $F$.

\begin{quote}
\textbf{Proposition} (page 23). \emph{Pour $F, G \in \mathfrak{F}$, sont
équivalents : $F \leq G$ ; $\widetilde{F} \subset \widetilde{G}$ ;
$\mathrm{Multomb}(F) \subset \mathrm{Multomb}(G)$ ; $\mathrm{Multomb}(F)$ est une
sous-figure de $\mathrm{Multomb}(G)$.}
\end{quote}

\emph{Corollaire.} Si $F \between G$, alors $\mathrm{Multomb}(F) \between
\mathrm{Multomb}(G)$, et
\[
\mathrm{Multomb}(F \vee G) = \mathrm{Multomb}(F) \cup \mathrm{Multomb}(G),
\]
\[
\mathrm{Multomb}(F \wedge G) = \mathrm{Multomb}(F) \cap \mathrm{Multomb}(G),
\]
parce que $\widetilde{F \vee G} = \widetilde{F} \cup \widetilde{G}$ et
$\widetilde{F \wedge G} = \widetilde{F} \cap \widetilde{G}$\footnote{La page
écrit « $F, G \in \mathcal{M}$ » pour $F, G \in \mathfrak{F}$. La première
égalité de déploiements vient de la forme « premier » (1.20) de
l'irréductibilité.}. La réciproque est fausse, même pour deux lieux distincts
$x$, $y$ : leurs multiombres $\lbrace \lbrace x \rbrace \rbrace$ et $\lbrace
\lbrace y \rbrace \rbrace$ sont disjointes, sans que $x$ et $y$ soient
nécessairement compatibles ; elle sera vraie pour les ateliers
ensemblistes\footnote{La page écrit « $\mathrm{Multomb}(X)$ »,
« $\mathrm{Multomb}(Y)$ » pour les figures réduites aux lieux $x$, $y$. En marge :
« vérifier plus tard ».}. Comme $\mathrm{Multomb}(\emptyset_{\mathfrak{F}}) =
\emptyset$ (1.40), deux figures disjointes ont des multiombres disjointes
(corollaire 2).

\pagerange{24}{28}
\subsection*{Le raffinement se lit sur les multiombres (pages 24 à 28)}

\begin{quote}
\textbf{Proposition} (pages 24 à 27). \emph{$G \ll F$ si et seulement si
$\mathrm{Multomb}(G) \ll \mathrm{Multomb}(F)$.}
\end{quote}

Par At~4, $\mathrm{Omb}(F)$ est la réunion $|\mathrm{Multomb}(F)|$ des
$\mathrm{Omb}(X)$, $X \in \widetilde{F}$, et l'application
\[
\varphi_{F} : \mathrm{Omb}(F) \longrightarrow \widetilde{F}
\]
qui envoie $Y$ sur le plus petit élément de $F^{Y}$ (At~5) est celle qui envoie
chaque point du support de $\mathrm{Multomb}(F)$ sur sa strate ouverte ; elle
induit l'identité sur $\widetilde{F}$. Dire que $\mathrm{Multomb}(G) \ll
\mathrm{Multomb}(F)$, c'est dire que $\mathrm{Omb}(G) \subset \mathrm{Omb}(F)$
et qu'il existe une application croissante $f$ rendant commutatif le carré
(1.42)\footnote{Le numéro du carré se lit plutôt « 1.48 » ; ses propres renvois
des pages 26 et 27 à « (1.42) » le fixent. De même le numéro (1.43) ci-dessous.}

\begin{tikzcd}
\mathrm{Omb}(G) \arrow[r, hook] \arrow[d, "\varphi_{G}"'] & \mathrm{Omb}(F) \arrow[d, "\varphi_{F}"] \\
\widetilde{G} \arrow[r, "f"] & \widetilde{F}
\end{tikzcd}

Si $G \ll F$, on prend pour $f(X)$ le plus petit élément de $F^{X}$ ; $f$ est
croissante. La commutativité de (1.42) est la propriété que la page 26 énonce en
corollaire et propose aussitôt de prendre pour axiome. Elle se dit avec le
\emph{raffinement intérieur} (1.43) :
\[
X \mathrel{\overset{\circ}{\ll}} Y \iff X \in \mathrm{Omb}(Y)^{\circ} =
\mathrm{Omb}(Y) \smallsetminus \bigcup_{Y' < Y} \mathrm{Omb}(Y') ,
\]
c'est-à-dire $X \ll Y$, et $X \ll Y' \leq Y$ entraîne $Y' = Y$. Pour $Y$ dans un
déploiement $\widetilde{F}$, qui est fermé, $Y$ est le plus petit élément de
$F^{X}$ si et seulement si $X \mathrel{\overset{\circ}{\ll}} Y$.

\begin{quote}
\textbf{At~7.} La relation $\overset{\circ}{\ll}$ est transitive.
\end{quote}

C'est exactement ce qui fait commuter (1.42) : pour $X \in \mathrm{Omb}(G)$, si
$X \mathrel{\overset{\circ}{\ll}} X' \in \widetilde{G}$ et $X'
\mathrel{\overset{\circ}{\ll}} X'' \in \widetilde{F}$, alors $X
\mathrel{\overset{\circ}{\ll}} X''$, et $\varphi_{F}(X) = X'' =
f(\varphi_{G}(X))$\footnote{La page 26 laisse la vérification en suspens (« il
faut vérifier la commutativité de (1.42), ce qui … », illisible) ; la
déduction par At~7 est la nôtre, et c'est celle qu'appelle la page 27, « On
doit donc poser At~7 ».}. Réciproquement, s'il existe un tel $f$, chaque $X \in
\widetilde{G}$ vérifie $X \ll \varphi_{F}(X) \in \widetilde{F}$, donc $G \ll F$
par At~4 ; la croissance de $f$ n'a pas servi. D'où le corollaire 2 de la page
27 : \emph{$G \ll F$ si et seulement si $\mathrm{Omb}(G) \subset \mathrm{Omb}(F)$
et la partition de $\mathrm{Omb}(G)$ par les fibres de $\varphi_{G}$ est plus
fine que celle que $\varphi_{F}$ y induit}\footnote{Le corollaire est écrit avec
les rôles de $F$ et $G$ échangés. Une marge exprime un doute (« Ce ne sera
… vrai, il faut … »), lue en partie. L'énoncé est juste tel quel : la
condition sur les fibres fournit $f$, d'où $G \ll F$ par la réciproque, et le
sens direct rend alors $f$ croissante.}.

\pagerange{28}{33}
\section*{Le théorème-scholie (pages 28 à 33)}

\emph{Scholie} (page 28). Par $F \mapsto \mathrm{Multomb}(F)$, $(\mathfrak{F},
\leq, \ll)$ s'identifie à un « sous-atelier » de $\mathrm{Fig}(\mathcal{M})$.
Se donner l'atelier revient donc à se donner un ensemble $\mathcal{M}$ et une
partie (1.44)
\[
\mathfrak{F}_{\mathcal{M}} \subset \mathrm{Fig}(\mathcal{M}) \subset
\mathfrak{P}(\mathfrak{P}(\mathcal{M}))
\]
soumise à des conditions qui traduisent les précédentes. Du coup, dit-il, la
multiombre reprend de l'intérêt et mériterait un nom plus lapidaire.

Soit $\mathcal{M}^{*} \subset \mathfrak{P}(\mathcal{M})$ l'ensemble des strates
de toutes les figures de $\mathfrak{F}_{\mathcal{M}}$. Si $\mathcal{M}^{*}$ est
une figure de support $\mathcal{M}$ qui sépare les points, $X \mapsto
\mathrm{Omb}(X) = \mathcal{M}_{\ll X}$ est une bijection de $\mathcal{M}$ sur
$\mathcal{M}^{*}$ qui change $\ll$ en l'inclusion, et les éléments de
$\mathcal{M}^{*}$ sont les parties fermées irréductibles de $(\mathcal{M}, \ll)$
qui ont un point générique\footnote{Pour la topologie d'Alexandrov de l'ordre
$\ll$, dont les fermés sont les parties fermées au sens des conventions : les
fermés irréductibles à point générique sont les $\mathcal{M}_{\ll X}$. Le
vocabulaire de point générique est le sien ; le nom de topologie d'Alexandrov
(P.~Alexandrov, 1937) est le nôtre.}. Chaque $\Phi \in \mathfrak{F}_{\mathcal{M}}$
est alors l'image d'une unique partie $F \subset \mathcal{M}$, d'où un ensemble
$\widetilde{\mathfrak{F}} \simeq \mathfrak{F}_{\mathcal{M}}$ de parties de
$\mathcal{M}$.

\begin{quote}
\textbf{Théorème-scholie} (pages 32 et 33). \emph{La catégorie des ateliers
$(\mathfrak{F}, \leq, \ll)$ satisfaisant At~1 à At~7 équivaut à celle des
couples $(\mathcal{M}, \mathfrak{F}_{\mathcal{M}})$,
$\mathfrak{F}_{\mathcal{M}} \subset \mathrm{Fig}(\mathcal{M})$, satisfaisant :}
\end{quote}
\begin{enumerate}
\item[At$^{*}$1] \emph{$\mathfrak{F}_{\mathcal{M}}$ contient la figure vide et
est stable par réunions majorées dans $\mathfrak{F}_{\mathcal{M}}$ ;}
\item[At$^{*}$2] \emph{si $A$ est une strate commune à $\Phi, \Psi \in
\mathfrak{F}_{\mathcal{M}}$, les multistrates $\Phi_{A}$ et $\Psi_{A}$ qu'elle
définit dans l'une et dans l'autre --- ses sous-strates --- sont les mêmes, et
appartiennent à $\mathfrak{F}_{\mathcal{M}}$ ;}
\item[At$^{*}$3] \emph{si $A \subset B$ sont deux strates, la multistrate de
$A$ raffine celle de $B$ ;}
\item[At$^{*}$4] \emph{si $\Phi_{A} \cup \Psi_{B} \in \mathfrak{F}_{\mathcal{M}}$
pour toutes strates $A$ de $\Phi$ et $B$ de $\Psi$, alors $\Phi \cup \Psi \in
\mathfrak{F}_{\mathcal{M}}$ ;}
\item[At$^{*}$5] \emph{$\mathcal{M}^{*}$ est une figure de support $\mathcal{M}$,
qui sépare les points de $\mathcal{M}$.}
\end{enumerate}

« Toute strate porte canoniquement une structure de multistrate. » Corollaire
(page 31) : si $\Phi \subset \Psi$ dans $\mathfrak{F}_{\mathcal{M}}$, $\Phi$ est
fermée dans $\Psi$, donc en est une sous-figure. Et l'on retrouve sur
$\mathcal{M} \simeq \mathcal{M}^{*}$ l'ordre $\leq$ : $X \leq Y$ si et seulement
si $\mathrm{Omb}(X)$ est une strate de la multistrate $\Phi_{\mathrm{Omb}(Y)} =
\mathrm{Multomb}(Y)$, les éléments de $\widetilde{\mathfrak{F}}$ étant alors des
parties fermées et les $\mathcal{M}_{\leq X}$ ses éléments strictement
irréductibles (page 32). Et il conclut :
« J'espère que je n'ai pas oublié d'axiome ! »\footnote{La liste a été écrite deux
fois. Celle des pages 29 à 31 commence par la séparation des points (At$^{*}$1)
et place la stabilité par réunions majorées en At$^{*}$3 ; celle des pages 32 et
33 est reproduite ici, dans l'ordre où une flèche de la page 33 replace At$^{*}$3.
L'équivalence n'est pas démontrée, ni la catégorie précisée --- l'adjectif
« (quasi-) » devant « ateliers » est lu avec doute --- et rien n'assure que la
liste rende At~6 et At~7 ; on ne l'affirme pas. La marge de la page 30 renvoie à
la première mouture, pages 54, 55 et 58 de celle-ci.}

\pagerange{34}{40}
\section*{Petites multiombres, exemples, sous-ateliers (pages 34 à 40)}

Le 20 juin, il passe aux \emph{petites} multiombres. Comme $\mathrm{Fig}(\mathcal{M})
\to \mathrm{Fig}(\mathcal{L})$ est croissante pour $\leq$ et $\ll$, $F \mapsto
\mathrm{multomb}(F)$ l'est aussi, commute aux bornes supérieures et inférieures
pour $\leq$, envoie sous-figures sur sous-figures --- mais n'est plus
nécessairement injective.

\emph{Exemple 1 A.} $E$ est un espace topologique ; une figure est un triplet
$(X, \Sigma, \Phi)$ formé d'une partie fermée $X$, d'une structure linéaire par
morceaux modérée $\Sigma$ sur $X$ et d'une figure ensembliste $\Phi$ de support
$X$ dont les strates sont des fermés linéaires par morceaux. Sous-figure et
raffinement demandent que $X'$ soit un sous-polyèdre de $X$ de structure
induite, et respectivement $\Phi' \leq \Phi$, $\Phi' \ll \Phi$. Les lieux sont
les points de $E$, et $(X, \Sigma, \Phi) \mapsto \Phi$ n'est pas injective dès
que $E$ contient un arc : deux structures linéaires par morceaux sur un même
arc donnent la même figure ensembliste.

\emph{Exemple 1 B.} Les multistrates sont des simplexes affines $(A, \sigma)$,
$A \subset E$ ; $\leq$ est « face de », $\ll$ est le raffinement des figures
ensemblistes $\mathfrak{L}_{\sigma}$ associées, et deux simplexes sont compatibles
si leur intersection est vide ou une face commune. Les figures sont les
ensembles finis --- ou localement finis --- de simplexes, fermés pour « face
de », deux à deux compatibles\footnote{La page écrit, dans la condition de
compatibilité, $(X, \sigma)$, $(X', \sigma')$ pour $(A, \sigma)$, $(A',
\sigma')$.}.

\emph{Exemple 2}, un « pseudodisque » dans un espace, n'est lisible qu'en partie :
pour en faire une figure, la seule partie contractile ne suffit pas, il faut
encore une structure au bord.

Dans les exemples 1, il est « illusoire » d'espérer que $\mathrm{multomb}$
réfléchisse $\leq$ et $\ll$ ; c'est même faux. D'où une première
\emph{définition} : un atelier est \emph{ensembliste} si $\mathfrak{F} \to
\mathrm{Fig}(\mathcal{L})$ est injective et réfléchit $\leq$, $\ll$ et
$\between$, ce qui l'identifie à un « sous-atelier » de
$\mathrm{Fig}(\mathcal{L})$. La page 52 en donnera la forme définitive.

\emph{Morphismes} (pages 36 et 37). Une application $f : \mathfrak{F} \to
\mathfrak{F}'$ est un \emph{morphisme d'ateliers} si elle est croissante pour
$\leq$ et $\ll$ et commute aux bornes supérieures majorées ; elle est
\emph{compatible aux strates} si elle envoie multistrates sur multistrates. Une
telle $f$ est déterminée par $f_{\mathrm{mult}} : \mathcal{M} \to \mathcal{M}'$,
qui doit être croissante pour $\leq$ et $\ll$ et envoyer les déploiements sur des
déploiements, puisque $f(F)$ est la borne supérieure des $f(X)$, $X \in
\widetilde{F}$\footnote{La page 37 énonce cette équivalence et un NB sur les
isomorphismes, écrit serré et lu en partie ; on ne les démontre pas.}.

\emph{Sous-ateliers et parties spécieuses} (pages 37 à 40). Après une rédaction
barrée (page 38), la page 39 reprend sur un ensemble ordonné $(\mathfrak{F}, \leq)$
avec bornes supérieures majorées et plus petit élément (At~1). Soit
$\mathfrak{F}'$ une partie fermée. Elle est non vide si et seulement si elle
contient $\emptyset_{\mathfrak{F}}$ ; elle est automatiquement stable par les
bornes supérieures des familles majorées \emph{dans $\mathfrak{F}'$}, qui sont les
mêmes dans $\mathfrak{F}$ et dans $\mathfrak{F}'$ ; ses éléments irréductibles
sont $\mathcal{M}' = \mathcal{M} \cap \mathfrak{F}'$, et elle hérite de At~2.
Elle est dite \emph{spécieuse} si elle est stable par les bornes supérieures des
familles majorées \emph{dans $\mathfrak{F}$}\footnote{Une définition barrée de la
page 39 appelait « ample » ce que la page 40 montre automatique ; la marge de la
page 40 garde le mot sans qu'on le lise assez pour l'employer.}.

L'exemple de la page 40 montre la différence. Pour $\mathfrak{F} =
\mathfrak{P}(I)$ ordonné par inclusion, les multistrates sont les singletons, et
une partie fermée non vide est une famille de parties de $I$ stable par passage
aux parties, c'est-à-dire un \emph{complexe simplicial abstrait} de sommets dans
$I$ --- si $I$ est fini, ou si l'on se borne aux simplexes finis. Elle est
spécieuse si la réunion de deux simplexes en est un, c'est-à-dire si elle est
$\mathfrak{P}(I')$ pour une partie $I'$ de $I$ : « c'est pas du tout
pareil ! »\footnote{Le nom de complexe simplicial abstrait est le sien.}

\pagerange{41}{46}
\section*{Sous-ateliers spécieux et axiome des lieux (pages 41 à 46)}

\begin{quote}
\textbf{Proposition} (page 41). \emph{Sous At~1 et At~2, les applications
\[
\mathfrak{F}' \longmapsto \mathfrak{F}' \cap \mathcal{M}, \qquad \mathcal{M}'
\longmapsto \mathrm{Env}(\mathcal{M}') = \lbrace F \in \mathfrak{F} \mid
\widetilde{F} \subset \mathcal{M}' \rbrace
\]
sont des bijections réciproques entre parties fermées spécieuses de
$\mathfrak{F}$ et parties fermées de $\mathcal{M}$.} Elles ne dépendent que de
$\leq$.
\end{quote}

C'est la forme, pour des familles de parties fermées, de la représentation de
Birkhoff\footnote{La page n'écrit pas la démonstration. $\mathrm{Env}(\mathcal{M}')$
est spécieuse parce que les strates d'une borne supérieure sont la réunion des
strates (forme « premier » de (1.20)) ; inversement, une figure dont les strates
sont dans une partie spécieuse $\mathfrak{F}'$ est la borne supérieure de ses
strates (At~2), famille majorée dans $\mathfrak{F}$, donc est dans
$\mathfrak{F}'$. Les hypothèses At~1, At~2 sont ajoutées.}.

\emph{Sous-ateliers.} Pour que $\mathfrak{F}'$, muni des ordres induits, soit un
atelier, At~1, At~2 sont acquis, At~4 à At~7 passent « trivialement » --- le
raffinement intérieur est le même dans $\mathfrak{F}$ et dans $\mathfrak{F}'$ ---,
mais At~3 demande une condition. D'où la \emph{définition} (page 42) : un
\emph{sous-atelier} de $\mathfrak{F}$ est une partie satisfaisant
\begin{enumerate}
\item[Sousat~1] \emph{$\mathfrak{F}'$ est une partie fermée non vide de
$\mathfrak{F}$ ;}
\item[Sousat~2] \emph{si $F, G \in \mathfrak{F}'$ et si, pour toutes strates $X$
de $F$ et $Y$ de $G$, $\lbrace X, Y \rbrace$ est majoré dans $\mathfrak{F}'$,
alors $F \vee G \in \mathfrak{F}'$.}
\end{enumerate}
Une partie fermée spécieuse est un sous-atelier, dit \emph{spécieux}. L'image de
$\mathrm{Multomb} : \mathfrak{F} \to \mathrm{Fig}(\mathcal{M})$ est un
sous-atelier, en général non spécieux. Les sous-ateliers d'un sous-atelier sont
les sous-ateliers contenus dans celui-ci, et de même pour les spécieux (1.50),
(1.51).

\emph{Ateliers triviaux} (page 44). Un ensemble ordonné $(\mathcal{M}, \leq)$, un
ensemble $\mathfrak{F}$ de parties fermées recouvrant $\mathcal{M}$ et stable par
passage aux parties fermées, et $\ll \overset{\mathrm{df}}{=} \leq$ : si At~3 est
satisfait, c'est un atelier au sens de At~1 à At~7. Les axiomes suivants,
pense-t-il, élimineront cet exemple.

\subsection*{L'axiome des lieux}

Pour $X \in \mathcal{M}$, soit $\partial X$ la borne supérieure de ses strates
distinctes de $X$.

\begin{quote}
\textbf{At~8} (page 45). \emph{Pour toute $X \in \mathcal{M}$, il existe un lieu
$x \ll X$ qui ne raffine pas $\partial X$, c'est-à-dire $x
\mathrel{\overset{\circ}{\ll}} X$.} Autrement dit, (1.52) :
\[
\boxed{\ \mathrm{omb}(X)^{\circ} = \mathrm{omb}(X) \smallsetminus
\mathrm{omb}(\partial X) = \mathrm{omb}(X) \smallsetminus \bigcup_{Y
\mathrel{\triangleleft} X,\ Y \neq X} \mathrm{omb}(Y) \neq \emptyset\ }
\]
\end{quote}

La première mouture énonçait cette propriété en proposition, sans démonstration ;
elle devient ici un axiome\footnote{Dossier 156-4, page 88, où les membres de la
multiombre sont dits avoir un intérieur non vide. La page 72 du présent dossier
cite cette formule sous le numéro « (1.48) » ; c'est (1.52).}. Alors
$\mathrm{multomb}(F)$ est indexée par $\widetilde{F}$ tout entier, et :

\emph{Corollaire} (pages 45 et 46). Pour $X, Y$ strates d'une même figure, sont
équivalents : $X \leq Y$ ; $\mathrm{omb}(X) \subset \mathrm{omb}(Y)$ ;
$\mathrm{omb}(X)^{\circ} \subset \mathrm{omb}(Y)$ ; $\mathrm{omb}(X)^{\circ} \cap
\mathrm{omb}(Y) \neq \emptyset$\footnote{Pour le dernier pas : si $x \in
\mathrm{omb}(X)^{\circ} \cap \mathrm{omb}(Y)$, le plus petit élément de
$F^{x}$ (At~5) est $\leq X$ et $\leq Y$, et ne peut être une strate stricte de
$X$ ; c'est donc $X$, et $X \leq Y$.}. D'où une bijection (1.53)
\[
\mathrm{Ssfig}(F) \xrightarrow{\ \sim\ } \mathrm{Ssfig}(\mathrm{multomb}(F)), \qquad
F' \longmapsto \mathrm{multomb}(F'),
\]
et $F$ est une multistrate si et seulement si $\mathrm{multomb}(F)$ est
élémentaire.

At~8 n'a aucune raison de passer à un sous-atelier, même spécieux. Il passe aux
sous-ateliers \emph{pleins}, fermés aussi pour $\ll$ ($F \in \mathfrak{F}'$, $G
\ll F \Rightarrow G \in \mathfrak{F}'$)\footnote{La page écrit « $G \ll F$
implique $G \ll F$ » ; on attend $G \in \mathfrak{F}'$.}, et il suffit qu'ils
soient \emph{$\mathcal{L}$-pleins} ($x \in \mathcal{L}$, $x \ll F \in
\mathfrak{F}' \Rightarrow x \in \mathfrak{F}'$). En général les lieux de
$\mathfrak{F}$ qui sont dans $\mathfrak{F}'$ sont des lieux de $\mathfrak{F}'$,
$\mathcal{L} \cap \mathfrak{F}' \subset \mathcal{L}'$ ; pour un sous-atelier
$\mathcal{L}$-plein il y a égalité, et ombres et multiombres sont les mêmes dans
$\mathfrak{F}$ et dans $\mathfrak{F}'$\footnote{L'égalité utilise At~8 dans
$\mathfrak{F}$ : une figure non vide de $\mathfrak{F}'$ est raffinée par un lieu
de $\mathfrak{F}$, qui est dans $\mathfrak{F}'$.}.

\pagerange{47}{52}
\section*{Quasi-morphismes et plongements (pages 47 à 52)}

\emph{Définition} (page 47). Un \emph{quasi-morphisme} $f : \mathfrak{F}' \to
\mathfrak{F}$ est une application croissante pour $\leq$ et $\ll$, qui commute
aux bornes supérieures majorées --- donc envoie la figure vide sur la figure
vide --- et aux bornes inférieures pour $\leq$. On réserve le nom de morphisme
pour plus tard. Si $f$ envoie multistrates sur multistrates, elle est connue par
sa restriction $f_{\mathrm{str}} : \mathcal{M}' \to \mathcal{M}$, et elle est
injective si et seulement si $f_{\mathrm{str}}$ l'est.

La condition qui compte est que, pour toute $F' \in \mathfrak{F}'$ et $F =
f(F')$, l'application
\[
(1.54) \qquad \mathrm{Ssfig}(F') \longrightarrow \mathrm{Ssfig}(F), \qquad G'
\longmapsto f(G')
\]
soit bijective\footnote{La page 48 écrit $\mathrm{Ssfig}(F')$ des deux côtés ;
« où $F = f(F')$ » fixe le second. Le numéro (1.54), cité ensuite, n'est pas
écrit à côté de la formule.}. Dans les cas qu'il connaît, (1.54) provient d'une
application croissante $\varphi : \widetilde{F} \to \widetilde{F'}$ (1.55) --- en
sens inverse de $f$ --- par $A \mapsto \varphi^{-1}(A)$ sur les parties fermées
(1.56), et elle est toujours au moins surjective. Une telle $\varphi$ est unique,
en vertu du

\begin{quote}
\textbf{Lemme} (page 49). \emph{Soient $I$, $I'$ deux ensembles ordonnés.
L'application $\varphi \mapsto \varphi^{-1}$,
\[
\mathrm{Hom}_{\mathrm{croiss}}(I, I') \longrightarrow
\mathrm{Hom}(\mathfrak{P}_{\mathrm{f}}(I'), \mathfrak{P}_{\mathrm{f}}(I)),
\]
est injective, et son image est formée des applications qui commutent aux
réunions et aux intersections quelconques.}
\end{quote}

La page est très retravaillée, et c'est l'énoncé exact que l'on
donne\footnote{La page parle d'ensembles « préordonnés » ; pour un préordre il
faut passer à l'ordre associé, sans quoi l'injectivité tombe. Elle décrit
l'application en mots dans le sens contraire de la flèche, et mentionne une
hypothèse de sobriété et le « sobrifié » $\widehat{I}'$ de $I'$. Si l'on ne
demande que la commutation aux intersections quelconques et aux réunions finies,
on obtient les applications continues de $I$ dans le sobrifié de $I'$, pour les
topologies d'Alexandrov : c'est sans doute ce que visait la mention des
morphismes de topos $\mathrm{Top}(I) \to \mathrm{Top}(I')$. Cette lecture est la
nôtre. L'énoncé exact du lemme est connu sous le nom de correspondance de
G.~N.~Raney (1952) ; la notion d'espace sobre est celle de SGA~4.}. La surjectivité
de (1.54) équivaut alors à ce que l'ordre de $\widetilde{F}$ soit l'image
réciproque par $\varphi$ de celui de $\widetilde{F'}$.

\emph{Que cette condition est substantielle} (pages 49 et 50). Dans les ateliers
triviaux de la page 44 avec $\mathfrak{F} = \mathfrak{P}_{\mathrm{f}}(\mathcal{M})$,
$\mathfrak{F}' = \mathfrak{P}_{\mathrm{f}}(\mathcal{M}')$, les quasi-morphismes
$\mathfrak{F}' \to \mathfrak{F}$ correspondent, par le lemme, à toutes les
applications croissantes $\mathcal{M} \to \mathcal{M}'$, sans aucune condition :
la surjectivité de (1.54) est donc une vraie exigence. En passant, la page
examine At~8 dans ces ateliers et la trouve « des plus restrictives » ; elle est
en fait plus forte encore\footnote{Avec $\ll = \leq$, les lieux sont les
éléments minimaux de $\mathcal{M}$, et $\mathrm{omb}(\partial X)$ contient tout
élément minimal $x < X$. La condition de la page, « $x$ minimal, $x \leq X$, et
pas de $Y$ avec $x \leq Y < X$ », exclut déjà $Y = x$ dès que $x \neq X$ : At~8
force tout élément de $\mathcal{M}$ à être minimal, c'est-à-dire $\mathcal{M}$ à
être discret.}.

\emph{Exemples} (pages 50 et 51). Une application d'ensembles $f : \mathcal{L}
\to \mathcal{L}'$ induit un quasi-morphisme
\[
f^{*} : \mathrm{Fig}(\mathcal{L}') \longrightarrow \mathrm{Fig}(\mathcal{L}),
\qquad \Phi' \longmapsto \lbrace f^{-1}(A') \mid A' \in \Phi',\ f^{-1}(A') \neq
\emptyset \rbrace ,
\]
contravariant en $f$. $\mathfrak{F} \to \mathrm{Fig}(\mathcal{M})$ est un
quasi-morphisme injectif qui satisfait (1.54) ; composé avec le $f^{*}$ de
$\mathcal{L} \hookrightarrow \mathcal{M}$, il donne $\mathfrak{F} \to
\mathrm{Fig}(\mathcal{L})$, $F \mapsto \mathrm{multomb}(F)$, qui satisfait (1.54)
sous At~8 --- c'est (1.53).

\emph{Plongements} (page 51). Un quasi-morphisme $f$ est un \emph{plongement} si
son image est un sous-atelier et si $f$ induit un isomorphisme sur elle. Il
l'affirme équivalent à
\begin{enumerate}
\item[a')] pour toute $F'$, $\mathrm{Sousfig}(F') \to \mathrm{Sousfig}(f(F'))$
est bijective --- la surjectivité suffit, vu b') ;
\item[b')] $f(G') \ll f(F')$ entraîne $G' \ll F'$.
\end{enumerate}
La démonstration de la page 52 établit ce qui en est l'essentiel : b') rend $f$
injective, puisque $\ll$ est antisymétrique, et a') avec b') rend $f$
réfléchissante pour $\leq$ ; que l'image satisfasse Sousat~2 n'est pas
vérifié\footnote{La page renvoie, pour l'injectivité, à « b) » ; c'est b'). Une
marge s'interroge sur l'accord avec la définition de la page 47.}.

\pagerange{52}{58}
\section*{Ateliers ensemblistes (pages 52 à 58)}

\begin{quote}
\textbf{Définition} (page 52). Un atelier satisfaisant At~1 à At~8 est
\emph{ensembliste} si le quasi-morphisme canonique $\mathfrak{F} \to
\mathrm{Fig}(\mathcal{L})$ est un plongement, ou, ce qui revient au même vu
(1.53) :
\end{quote}
\begin{enumerate}
\item[At ens~1] \emph{$\mathrm{multomb}(F) \ll \mathrm{multomb}(G) \Rightarrow F
\ll G$,}
\end{enumerate}
qu'il suffit de demander pour des multistrates.

On n'exige pas que le sous-atelier image soit spécieux. Exemple : $E$ un espace
modéré, les multistrates ses simplexes topologiques ou modérés, sans structure
affine cette fois, la compatibilité étant celle de l'exemple 1 B ; At~8 est
évident, et $\mathcal{L} = E$. La compatibilité des simplexes y est plus fine
que celle des figures ensemblistes associées, puisqu'on exige que
l'intersection soit une multistrate commune\footnote{La marge de la page 52
renvoie à « GF VI p.~12 », c'est-à-dire à la page 12 de ce dossier.}.

\begin{quote}
\textbf{Scholie} (pages 54 et 55). \emph{Se donner un atelier (quasi-)ensembliste
revient à se donner un ensemble $\mathcal{L}$ et un sous-atelier $\mathfrak{F}$
de $\mathrm{Fig}(\mathcal{L})$ tel que}
\end{quote}
\begin{enumerate}
\item[Atens$^{*}$1] \emph{pour tout $x \in \mathcal{L}$, $\lbrace \lbrace x
\rbrace \rbrace \in \mathfrak{F}$,}
\end{enumerate}
\emph{les conditions de sous-atelier s'écrivant}
\begin{enumerate}
\item[Atens$^{*}$2] \emph{$\mathfrak{F}$ est fermé dans $\mathrm{Fig}(\mathcal{L})$
pour $\leq$ ;}
\item[Atens$^{*}$3] \emph{si toute multistrate de $F$ et toute multistrate de $G$
sont majorées ensemble dans $\mathfrak{F}$, alors $F \vee G \in \mathfrak{F}$.}
\end{enumerate}

$\mathcal{L}$ se retrouve comme l'ensemble des lieux : $x \mapsto \lbrace \lbrace
x \rbrace \rbrace$ est une bijection sur les éléments minimaux de $\mathfrak{F}
\smallsetminus \lbrace \emptyset \rbrace$ pour $\ll$, car toute figure non vide
$F$ est raffinée par les $\lbrace \lbrace x \rbrace \rbrace$, $x \in
|F|$\footnote{C'est la définition des lieux de la première mouture (dossier
156-4, page 87). Une première forme de Atens$^{*}$1, que les strates séparent les
points, est barrée.}.

On peut aussi partir d'un ensemble $\mathcal{M}$ de figures ensemblistes
élémentaires dans $\mathcal{L}$ et d'une relation réflexive et symétrique
$\between_{\mathcal{M}}$, la « compatibilité géométrique », soumis à
\begin{enumerate}
\item[Atens$^{\sharp}$1] \emph{$\lbrace \lbrace x \rbrace \rbrace \in \mathcal{M}$
pour tout $x \in \mathcal{L}$ ;}
\item[Atens$^{\sharp}$2] \emph{toute sous-multistrate d'un élément de
$\mathcal{M}$ est dans $\mathcal{M}$ ;}
\item[Atens$^{\sharp}$3] \emph{$\between_{\mathcal{M}}$ est stable par passage aux
sous-multistrates, et entraîne la compatibilité ensembliste.}
\end{enumerate}
\emph{Scholie} (page 56) : les ateliers ensemblistes à figures de type fini
correspondent aux données $(\mathcal{L}, \mathcal{M}, \between_{\mathcal{M}})$
satisfaisant Atens$^{\sharp}$1 à Atens$^{\sharp}$3, toute sous-figure d'un $X \in
\mathcal{M}$ étant de type fini ; $\mathfrak{F}$ est l'ensemble des figures
ensemblistes de type fini dont les multistrates sont dans $\mathcal{M}$ et deux
à deux compatibles pour $\between_{\mathcal{M}}$\footnote{La page omet la
condition de compatibilité deux à deux, sans laquelle $\between_{\mathcal{M}}$ ne
jouerait aucun rôle ; on l'ajoute. Elle devient automatique dans le corollaire
qui suit.}.

On ne doit pas s'attendre à ce que $\mathfrak{F}$ soit spécieux au sens fort :
la figure formée de tous les points de $\mathcal{L}$ n'a aucune raison d'y être.
La condition raisonnable est que $\mathfrak{F}$ soit \emph{finiment spécieux},
stable par bornes supérieures \emph{finies} majorées dans
$\mathrm{Fig}(\mathcal{L})$ ; vu Atens$^{\sharp}$3, cela signifie que
$\between_{\mathcal{M}}$ est la compatibilité ensembliste. \emph{Corollaire} :
les ateliers ensemblistes finiment spécieux de type fini correspondent aux
couples $(\mathcal{L}, \mathcal{M})$ satisfaisant Atens$^{\sharp}$1 et
Atens$^{\sharp}$2, les sous-figures des $X \in \mathcal{M}$ étant de type
fini\footnote{La page écrit « sous-figure de $\mathcal{L}$ » pour sous-figure de
$X$.}.

\emph{Remarque} (page 58). L'autre cas intéressant qu'il ait rencontré est celui
où $X \between_{\mathcal{M}} Y$ signifie : $X$ et $Y$ sont ensemblistement
compatibles et $X \cap Y$ est vide ou est une multistrate commune --- le cas
simplicial. Alors (1.57) : pour toute figure $F$, deux strates de $F$ qui se
rencontrent ont pour borne inférieure leur intersection, qui est une
multistrate\footnote{La page dit que $\widetilde{F}$ admet les bornes
inférieures de deux éléments ; c'est faux pour deux strates disjointes, puisque
la figure vide n'est pas une multistrate. Ce qui est vrai : $\widetilde{F} \cup
\lbrace \emptyset \rbrace$ a des bornes inférieures binaires.}.

\pagerange{59}{64}
\section*{Axiomes de disjonction (pages 59 à 64)}

At~8 était « l'axiome des lieux ». Viennent les \emph{axiomes de disjonction}.

\emph{Proposition} (page 59). Sous At~1 à At~3, $F \parallel G$ si et seulement
si toute strate de $F$ est disjointe de toute strate de $G$. On voudrait de plus
la stabilité de la disjonction par raffinement :
\[
(1.58) \qquad F \parallel G, \quad F' \ll F, \quad G' \ll G \ \Longrightarrow\ F'
\parallel G' ,
\]
qu'on ramène, par At~4, au cas de multistrates. Elle ne résulte pas des
conditions Atens$^{\sharp}$, même pour un atelier ensembliste de type fini. Le
contre-exemple de la page 60 : trois multistrates à une seule strate, $X =
\lbrace A \rbrace$, $Y = \lbrace B \rbrace$, $X' = \lbrace A' \rbrace$, avec $A'
\subsetneq A$ et $A \cap B = \emptyset$. On a $X' \ll X$, et la compatibilité
ensembliste des trois est acquise ; Atens$^{\sharp}$3 laisse donc libre de choisir
$X \between_{\mathcal{M}} Y$ et non $X' \between_{\mathcal{M}} Y$. Alors $X
\parallel Y$ mais $X' \nparallel Y$. D'où la condition :
\begin{enumerate}
\item[(*)] \emph{si $X' \ll X$, $Y' \ll Y$ dans $\mathcal{M}$ et $X \parallel Y$,
alors $X' \parallel Y'$,}
\end{enumerate}
qu'il ne veut pas poser en axiome : il y voit un cas particulier de l'axiome de
recollement des raffinements, At~11, qui viendra plus loin\footnote{La
parenthèse « il suffit de poser $X \between Y$ » qui suit (*) se comprend sans
doute ainsi : dans le cas ensembliste, l'intersection des supports de $X'$ et
$Y'$ est vide, et seule la compatibilité de $X'$ et $Y'$ est en question. C'est
notre lecture. Le nom de la condition, devant « Si », est fortement biffé.}.

\begin{quote}
\textbf{At~9} (page 61, « critère de disjonction par les lieux »). \emph{Si $X,
Y \in \mathcal{M}$ et si tout lieu de $X$ est disjoint de tout lieu de $Y$,
alors $X \parallel Y$.}
\end{quote}
\begin{quote}
\textbf{Atens~2} (pour les ateliers ensemblistes). \emph{Deux lieux distincts sont
compatibles} ; en termes de $\mathcal{M} \subset \mathrm{Fig}(\mathcal{L})$,
$\lbrace \lbrace x \rbrace \rbrace \between_{\mathcal{M}} \lbrace \lbrace y
\rbrace \rbrace$ pour $x \neq y$ (Atens$^{*}$4).
\end{quote}

Pour deux lieux, (1.59) :
\[
x \parallel y \iff x \between y \text{ et } x \neq y .
\]
\emph{Corollaire} (page 62). Dans un atelier ensembliste satisfaisant At~9 et
Atens~2, deux figures sont disjointes si et seulement si leurs lieux le sont deux
à deux --- c'est l'annonce (1.13) de la page 8\footnote{La page écrit
« $\mathrm{omb}\,X$, $\mathrm{omb}\,Y$ » pour $\mathrm{omb}\,F$,
$\mathrm{omb}\,G$. Le sens « figures disjointes $\Rightarrow$ lieux disjoints »
vient de ce que, dans le cas ensembliste, deux figures disjointes ont des
supports disjoints, et que deux lieux distincts sont disjoints par Atens~2 et
(1.59).}. Un second corollaire affirme (1.58) pour tout atelier satisfaisant
At~9\footnote{Pour cela il faudrait savoir que les lieux de deux multistrates
disjointes sont deux à deux disjoints, ce que la page ne montre pas à cet
endroit : c'est le cas ensembliste du corollaire précédent. La déduction de
(1.58) à partir de l'axiome de recollement, page 70, est en revanche complète ;
c'est à elle qu'on se fie.}.

\emph{Critères de compatibilité par les lieux.} Pour $X, Y \in \mathcal{M}$, soit
$F = X \cap Y = X \wedge Y$.

\begin{quote}
\textbf{At~9'} (page 62). \emph{Si tout lieu de $\mathrm{omb}(X) \smallsetminus
\mathrm{omb}(F)$ est disjoint de tout lieu de $\mathrm{omb}(Y) \smallsetminus
\mathrm{omb}(F)$, alors $X \between Y$.}
\end{quote}

L'hypothèse entraîne $\mathrm{omb}(X) \cap \mathrm{omb}(Y) = \mathrm{omb}(F)$,
puisqu'un lieu commun hors de $\mathrm{omb}(F)$ serait disjoint de lui-même ; la
réciproque, dit la marge, résultera d'un axiome de disjonction (page 72). Cet
axiome ferait reposer toute la description sur la seule relation $\between$ entre
lieux ; il est raisonnable, mais il n'est pas satisfait dans les cas simpliciaux
les plus importants, où l'on demande que l'intersection soit une multistrate.
D'où la variante « simpliciale » :

\begin{quote}
\textbf{At~9''} (page 63). \emph{$X \between Y$ si et seulement si tout lieu de
$\mathrm{omb}(X) \smallsetminus \mathrm{omb}(F)$ est disjoint de tout lieu de
$\mathrm{omb}(Y) \smallsetminus \mathrm{omb}(F)$, et si $F$ est vide ou est une
multistrate.}
\end{quote}

Chacun de At~9', At~9'' entraîne At~9\footnote{La page écrit seulement « $F$ est
élémentaire ». Sans le cas $F = \emptyset$, At~9'' rendrait incompatibles deux
simplexes disjoints, et n'entraînerait pas At~9 ; on l'ajoute, comme le fait la
page 58 pour la compatibilité simpliciale. L'implication utilise At~8 : si les
lieux de $X$ et de $Y$ sont deux à deux disjoints, $F$ n'a pas de lieu, donc est
vide.}.

\emph{At~10.} Une longue note oblique en tête de la page 64, écrite par-dessus le
texte, énonce un « axiome de disjonction » At~10, lisible par fragments seulement.
D'après ces fragments et l'usage qu'en font les pages 72 et 83, c'est le
\emph{critère de disjonction pour les lieux} :
\begin{quote}
\emph{si $X \neq Y$ sont deux multistrates compatibles, tout lieu de
$\mathrm{omb}(X)^{\circ}$ est disjoint de tout lieu de
$\mathrm{omb}(Y)^{\circ}$}\footnote{Restitution nôtre. Se lisent : « At~10 Soient
$X, Y \in \mathcal{M}$ … $X \neq Y$ … $x \in \mathrm{omb}(X)^{\circ}$ et
$y \in \mathrm{omb}(Y)^{\circ}$ … $x \parallel y$ ». La page 83 l'emploie
exactement sous cette forme, et c'est l'analogue de la condition Cont~7 de la
première mouture (dossier 156-4, pages 41 et suivantes : deux multistrates
compatibles distinctes ont des ombres ouvertes disjointes, lieu par lieu). La
note groupe At~8, At~9, At~10 sous le nom d'« axiomes des lieux ».}.
\end{quote}

\pagerange{64}{72}
\section*{Restreindre et recoller les raffinements (pages 64 à 72)}

\begin{quote}
\textbf{Proposition} (page 64, 21 juin). \emph{Soient $F' \ll F$ et $G \leq F$.
Alors $G' = \mathrm{Inf}^{\ll}(F', G)$ existe ; c'est la sous-figure de $F'$
définie par}
\[
(1.60) \qquad \widetilde{G'} = \lbrace X \in \widetilde{F'} \mid X \ll G \rbrace .
\]
\end{quote}

\begin{tikzcd}
F' \arrow[r, no head, "\ll" description] & F \\
G' \arrow[u, no head, "\leq" description] \arrow[r, no head, "\ll" description] & G \arrow[u, no head, "\leq" description]
\end{tikzcd}

La partie (1.60) est fermée dans $\widetilde{F'}$, donc définit une sous-figure
$G'$, et $G' \ll G$ par At~4. Pour voir que c'est la borne inférieure pour
$\ll$, il suffit par At~4 de traiter une multistrate $Z \ll F'$, $Z \ll G$, et,
$X'$ désignant la strate de $F'$ telle que $Z \mathrel{\overset{\circ}{\ll}} X'$,
de montrer $X' \ll G$ :

\emph{Lemme} (page 65). Si $F' \ll F$, $X' \leq F'$, $G \leq F$, $Z
\mathrel{\overset{\circ}{\ll}} X'$ et $Z \ll G$, alors $X' \ll G$. --- Soit $X$
la strate de $F$ telle que $X' \mathrel{\overset{\circ}{\ll}} X$. Par At~7, $Z
\mathrel{\overset{\circ}{\ll}} X$, donc $X$ est le plus petit élément de
$F^{Z}$. Par At~4, $Z \ll Y$ pour une strate $Y$ de $G$, qui est dans $F^{Z}$ ;
donc $X \leq Y$, et $X' \ll X \leq Y \leq G$.

On note $G' = F'|G$ --- la page écrit aussi $F'.G$ --- le raffinement de $G$
\emph{induit} par $F'$.

\emph{Corollaire 1.} Si $F' \leq F$, alors $F'|G = F' \wedge G$.

\emph{Corollaire 2.} Pour $F' \ll F$, l'application
\[
(1.61) \qquad \mathrm{Sousfig}(F) \longrightarrow \mathrm{Sousfig}(F'), \qquad G
\longmapsto F'|G ,
\]
commute aux bornes supérieures et inférieures : en termes de l'application
croissante $\varphi : \widetilde{F'} \to \widetilde{F}$ (qui envoie $X$ sur la
strate de $F$ qu'il raffine intérieurement), c'est $\varphi^{*} : A \mapsto
\varphi^{-1}(A)$ sur les parties fermées.

\emph{Corollaire 3.} Si $F = \mathrm{Sup}\,F_{i}$ ($F_{i} \leq F$), et $F'_{i} =
F'|F_{i}$, alors chaque $F'_{i}$ raffine $F_{i}$, $F'_{i}$ et $F'_{j}$ induisent
le même raffinement de $F_{i} \cap F_{j}$, et $F' = \mathrm{Sup}\,F'_{i}$. Un
raffinement de $F$ est donc connu par une « famille compatible » de raffinements
des $F_{i}$. Reste à savoir si toute famille compatible provient d'un
raffinement :

\begin{quote}
\textbf{At~11} (page 67, \emph{axiome de recollement des raffinements}).
\emph{Soient $F \in \mathfrak{F}$ et, pour toute strate $X$ de $F$, un raffinement
$F'_{X} \ll X$, tels que $F'_{Y} = F'_{X}|Y$ pour $Y < X$. Il existe alors $F' \ll
F$ tel que $F'|X = F'_{X}$ pour toute strate $X$.}
\end{quote}

$F'$ est unique par le corollaire 3. En termes actuels, les raffinements forment
un préfoncteur $G \mapsto \mathrm{Raff}(G)$ contravariant sur $(\mathfrak{F},
\leq)$, par restriction, et At~11 dit que $\mathrm{Raff}(F)$ est la limite des
$\mathrm{Raff}(X)$ sur les strates de $F$ : c'est une condition de
faisceau\footnote{Le rapprochement est le nôtre.}. Le contenu essentiel, note la
page 68, est de comparer $F'_{X}$ et $F'_{Y}$ quand $X$ et $Y$ ne sont pas
comparables ; son « noyau » est

\begin{quote}
\textbf{At~11'} (page 68). \emph{Si $X \between Y$ dans $\mathcal{M}$, $X' \ll
X$, $Y' \ll Y$ et $X'|T = Y'|T$ pour $T = X \cap Y$, alors $X' \between Y'$.}
\end{quote}

At~11 entraîne At~11' (on recolle $X'$ et $Y'$ sur les strates de $X \vee Y$),
et At~11' entraîne At~11 quand $\widetilde{F}$ est fini, ou plus généralement
quand $\mathrm{Sup}\,F'_{X}$ existe, avec $\widetilde{F'} = \bigcup
\widetilde{F'_{X}}$\footnote{La démonstration des pages 68 et 69 est en partie
barrée et lue en partie. La page 68 écrit « $F_{X}$ » pour $F'_{X}$.}. Corollaire
2 (page 69) : un système compatible de raffinements des $F_{i}$, $F = \mathrm{Sup}
\,F_{i}$, provient d'un unique raffinement de $F$ ; pour $I$ fini, At~11' suffit.

\emph{Sommes disjointes} (page 70). Si $F \parallel G$, les strates de $F \amalg
G = F \vee G$ sont celles de $F$ et celles de $G$, jamais comparables entre
elles ; At~11 donne donc une bijection
\[
\mathrm{Raff}(F) \times \mathrm{Raff}(G) \xrightarrow{\ \sim\ } \mathrm{Raff}(F
\amalg G), \qquad (F', G') \longmapsto F' \amalg G', \qquad H' \longmapsto (H'|F,
H'|G) ,
\]
et en particulier $F' \parallel G'$, puisque $F' \wedge G' = H'|(F \wedge G) =
H'|\emptyset_{\mathfrak{F}}$ est vide par (1.61) : c'est (1.58), maintenant
démontrée\footnote{Ce corollaire porte sur la page le numéro 3, déjà employé page
66.}.

\emph{Figures discrètes} (pages 71 et 72). Une figure est \emph{discrète} si
toutes ses strates sont des lieux. Elle est connue par la partie $A \subset
\mathcal{L}$ de ses strates, formée de lieux deux à deux disjoints, et une
partie finie de $\mathcal{L}$ est une figure discrète si et seulement si ses
lieux sont deux à deux disjoints ; l'ordre y étant discret, ses sous-figures
sont ses parties quelconques. \emph{Corollaire 4} : une partie $A$ définit une
figure discrète qui raffine $F$ si et seulement si $A \subset \mathrm{omb}(F)$ et
$A \cap \mathrm{omb}(X)$ est une figure pour toute strate $X$ de $F$ (par
exemple, finie et formée de lieux deux à deux disjoints)\footnote{Énoncé sans
démonstration. Suit un paragraphe lu en partie, qui rapproche ce corollaire de la
réciproque de At~9' et de At~10.}.

\emph{Corollaire 5} (page 72). Si $G$ est une sous-figure de $F$ distincte de $F$,
il existe un lieu $x \ll F$ disjoint de $G$. --- On prend une strate $X$ de $F$
qui n'est pas dans $G$ et $x \in \mathrm{omb}(X)^{\circ}$ (At~8) ; $x$ n'est pas
une sous-figure de $G$, par le corollaire de la page 45, et sa compatibilité avec
$G$ s'obtient par At~9 et At~10, ou par le recollement des raffinements.

\pagerange{73}{80}
\section*{Compatibilité des raffinements, $\Sigma$-ateliers (pages 73 à 80)}

Il veut renforcer At~11' :

\begin{quote}
\textbf{Proposition (?)} (page 73). \emph{Soient $F \between G$, $K = F \cap G$,
$F' \ll F$, $G' \ll G$. Pour que $F' \between G'$, il faut et il suffit que
$F'|K \between G'|K$ ; et alors $F' \vee G' \ll F \vee G$.}
\end{quote}

Si $F'$ et $G'$ sont compatibles, $H' = F' \vee G' \ll H = F \vee G$ par At~4, et
$K' = F' \wedge G' \ll K$ ; mieux,
\[
K' = \mathrm{Inf}^{\ll}(F', G') = \mathrm{Inf}^{\ll}(F'|K, G'|K) = F'|K \cap
G'|K ,
\]
deux sous-figures de $H'$, qui sont donc compatibles. La réciproque se lit sur un
diagramme dont le cœur est le cube suivant, où les flèches à crochet sont des
$\leq$ et les flèches verticales des $\ll$ ; la marge dit les carrés horizontaux
cartésiens et cocartésiens pour $\leq$ :

\begin{tikzcd}[column sep=small, row sep=small]
H'_{K} \arrow[rr, hook] \arrow[dr, hook] \arrow[dd, "\ll"'] & & H'_{F} \arrow[dr, hook] \arrow[dd] & \\
 & H'_{G} \arrow[rr, hook] \arrow[dd] & & H' \arrow[dd, "\ll"] \\
K \arrow[rr, hook] \arrow[dr, hook] & & F \arrow[dr, hook] & \\
 & G \arrow[rr, hook] & & H
\end{tikzcd}

Sous le carré du haut, le diagramme de la page 74 porte encore $F' \leq H'_{F}$,
$G' \leq H'_{G}$, $F'_{K} = F'|K$ et $G'_{K} = G'|K$, sous-figures de $F'$, $G'$
et de $H'_{K}$, et $K' \leq F'_{K}, G'_{K}$\footnote{La page 73 dessine une
première version, barrée ; la page 74 la reprend. Le dessin ci-dessus n'en garde
que les huit sommets qui forment le cube, les autres étant donnés en prose ; il
normalise les directions. La page 73 écrit « $K''$ » pour $K'$.}. Si l'on savait
$F'$ et $H'_{K} = F'_{K} \vee G'_{K}$ compatibles, on formerait $H'_{F}$ et
$H'_{G}$, qui induisent tous deux $H'_{K}$ sur $K$, et l'on recollerait. On est
ainsi ramené au

\emph{Cor.\ Main} (page 76). Si $G \trianglelefteq F$, $F' \ll F$ induit $F'_{G}$
sur $G$, et $F'_{G} \leq G'$ pour un raffinement $G'$ de $G$, alors $F'$ et $G'$
sont compatibles, et $F' \cap G' = F'_{G}$\footnote{« Main » est en anglais sur la
page, souligné avec « Cor. ».}.

Il le récrit pour des multistrates, $X \supseteq Y$, $X'
\mathrel{\overset{\circ}{\ll}} X$, $Y' \mathrel{\overset{\circ}{\ll}} Y$ et $X'_{Y}
= X'|Y$ : si $X'_{Y} \between Y'$, alors $X' \between Y'$ ; c'est, dit-il, un cas
particulier de At~11 ou de At~11', qui résulterait aussi de At~9'' et At~10, ou de
At~9' et At~10. Plutôt que de le déduire, il pose, sous une forme qui contient la
proposition :

\begin{quote}
\textbf{At~12} (pages 76 et 77, \emph{critère de compatibilité par raffinements}).
\emph{Soient $X' \ll X$ et $Y' \ll Y$ dans $\mathcal{M}$, avec $X \between Y$, et
$K = X \cap Y$. Si $X'|K \between Y'|K$, alors $X' \between Y'$.}
\end{quote}

Selon la page 77, At~12 entraîne At~11' (si $X'|K = Y'|K$, la compatibilité est triviale), At~10,
la proposition de la page 73, donc At~11 pour $\widetilde{F}$ fini et le
corollaire 2 de la page 69 pour les familles finies.

\emph{$\Sigma$-ateliers} (pages 77 et 78). Pour donner une forme commune à At~9'
et At~9'', soit $\Sigma$ une sous-catégorie pleine de celle des ensembles
ordonnés. Un atelier est un \emph{$\Sigma$-atelier} si chaque $\widetilde{X}$, $X
\in \mathcal{M}$, est dans $\Sigma$, et si l'on a

\begin{quote}
\textbf{At~9$_{\Sigma}$} (\emph{critère de compatibilité par les lieux}).
\emph{Pour $X, Y \in \mathcal{M}$ et $K = X \cap Y$, de sorte que $\widetilde{K} =
\widetilde{X} \cap \widetilde{Y}$, on a $X \between Y$ si et seulement si (i) tout
lieu de $\mathrm{omb}(X) \smallsetminus \mathrm{omb}(K)$ est disjoint de tout lieu
de $\mathrm{omb}(Y) \smallsetminus \mathrm{omb}(K)$, et (ii) $\widetilde{K} \in
\Sigma$.}
\end{quote}

Corollaire 1, « ancien At~9'' » : (i) et $K \in \mathcal{M}$ entraînent $X \between
Y$. Corollaire 2, « ancien At~9' » : si $\Sigma$ contient tous les
$\widetilde{F}$, $X \between Y$ équivaut à (i)\footnote{Pour retrouver At~9
(lieux deux à deux disjoints, donc $K$ vide), il faut que l'ensemble ordonné vide
soit dans $\Sigma$, ce que la page ne dit pas ; on le suppose. Dans (i), un signe
de compatibilité barré remplace la disjonction ; on garde la disjonction, qui
exclut $x = y$. Les dernières lignes de la page 78 sont barrées.}. Tous les
ateliers auxquels il a songé sont des $\Sigma$-ateliers pour $\Sigma$ convenable.

\emph{Le tableau des implications} (pages 79 et 80), dans sa version au propre :

\begin{tikzcd}[column sep=small]
 & \mathrm{At}\,9_{\Sigma} + \mathrm{At}\,10 \arrow[dl, Rightarrow] \arrow[dr, Rightarrow] & & \\
\mathrm{At}\,9_{\Sigma} \arrow[d, Rightarrow] & & \mathrm{At}\,12 \arrow[dl, Rightarrow] \arrow[d, Rightarrow] & \mathrm{At}\,11 \arrow[dl, Rightarrow] \\
\mathrm{At}\,9 & \mathrm{At}\,10 & \mathrm{At}\,11' &
\end{tikzcd}

avec, en légende, At~11 équivalent au recollement des raffinements dans $F =
\mathrm{Sup}\,F_{i}$ quelconque, At~11' à ce recollement pour $I$
fini\footnote{Un premier schéma, barré, précède celui-ci ; le premier tableau de
la page 79 donne aussi At~9$_{\Sigma} \Rightarrow$ At~12. On a reproduit les
flèches de la version au propre telles qu'elles sont décrites, sans en ajouter.}.
Le seul axiome qui échappe à At~9$_{\Sigma}$ + At~10 est At~11, parce qu'il
porte sur des bornes supérieures infinies. On le remplace donc par

\begin{quote}
\textbf{At~11 bis} (page 80). \emph{Si $(F'_{i})$ est une famille filtrante
croissante de raffinements de $F$ telle que, pour toute strate $X$ de $F$, les
$F'_{i}|X$ aient une borne supérieure, alors $\mathrm{Sup}\,F'_{i}$ existe,}
\end{quote}
qui résulte de At~11 et l'entraîne moyennant At~11' --- un axiome, note la
marge, qui n'a rien à voir avec les lieux. D'où deux systèmes « réduits » :
At~9$_{\Sigma}$, At~10, At~11 bis ; ou, si l'on veut éviter le choix de
$\Sigma$, At~9, At~11 bis, At~12 --- au prix du critère universel de
compatibilité par les lieux, « bien sympathique ». La page finit sur la
\emph{saturation} d'un atelier, dans la définition de laquelle il faudra inclure
At~9' = At~9$_{\overline{\Sigma}}$, $\overline{\Sigma}$ désignant tous les
ensembles ordonnés.

\pagerange{81}{88}
\section*{Subdivisions (pages 81 à 88)}

« J'en viens aux subdivisions, dont il n'avait pas encore été question jusqu'à
présent ! » Dans cette mouture, la subdivision n'est pas primitive : elle va
être \emph{définie} à partir du raffinement et de la disjonction, comme la page 12
l'avait entrevu en (1.25)\footnote{Dans la première mouture, $\preccurlyeq$ est
une relation primitive dont se déduit $\ll$ (dossier 156-4, pages 52 à 60 et 84),
et le recollement des subdivisions y est l'axiome C8 (page 86 de ce
dossier-là).}.

\emph{Proposition} (page 81). Si $F' \ll F$ et s'il existe une figure non vide
$G$ disjointe de $F'$ mais non de $F$, on peut en trouver une avec de plus $G \ll
F$. Pour la démontrer, il isole :

\emph{Lemme} (page 82). Soient $F' \ll F$, $G \leq F$, $F'_{G} = F'|G$ et $H \ll
G$. Alors $H \parallel F'$ si et seulement si $H \parallel F'_{G}$. --- Le sens
direct est clair, $F'_{G}$ étant une sous-figure de $F'$. Pour l'autre, par At~9
il suffit de voir qu'un lieu $x$ de $H$ et un lieu $x'$ de $F'$ sont disjoints.
Si $x'$ raffine $F'_{G}$, c'est l'hypothèse. Sinon $x'$ ne raffine pas $G$ (il
raffinerait $\mathrm{Inf}^{\ll}(F', G) = F'_{G}$), donc $x$ et $x'$ sont dans
les ombres ouvertes de deux strates différentes de $F$, et $x \parallel x'$ par
At~10.

\begin{quote}
\textbf{At~13} (page 83). \emph{Soient $X \in \mathcal{M}$ et $F \ll X$. S'il existe
une figure non vide $G$ avec $G \parallel F$ et $G \nparallel X$ --- on peut alors
la prendre dans $\mathcal{L}$ ---, il existe un lieu $x$ de $X$ disjoint de $F$.}
\end{quote}

La marge range At~8, At~9 et At~13 parmi les axiomes « beaucoup de lieux ».

\begin{quote}
\textbf{Corollaire} (pages 83 et 84). \emph{Pour $F' \ll F$, sont équivalents :}
\end{quote}
\begin{enumerate}
\item[(i)] \emph{toute figure disjointe de $F'$ est disjointe de $F$ ;}
\item[(ii)] \emph{aucun lieu de $F$ n'est disjoint de $F'$ ;}
\item[(iii)] \emph{$F'$ est maximal, pour $\leq$, parmi les raffinements de $F$.}
\end{enumerate}

(i) $\Leftrightarrow$ (ii) est la proposition, sous At~13. Si (ii) est en défaut,
avec $x \in \mathrm{omb}(F)$ disjoint de $F'$, alors $F' \vee x$ est un
raffinement de $F$ (At~4) qui majore strictement $F'$, et (iii) est en défaut.
Si (iii) est en défaut, avec $F' < \overline{F'} \ll F$, une strate $X$ de
$\overline{F'}$ hors de $F'$ et $x \in \mathrm{omb}(X)^{\circ}$ donnent un lieu de
$F$ disjoint de $F'$\footnote{La page 83 écrit (ii) avec « $\mathrm{omb}(X)$ » ;
la démonstration de la page 84 se sert de $\mathrm{omb}(F)$. Elle écrit aussi
« $x \parallel F$ » pour $x \parallel F'$, et un $\leq$ barré pour $\neq$. Que $x$
soit disjoint de $F'$ s'obtient comme au corollaire 5 de la page 72.}. Une
quatrième condition, la maximalité pour $\ll$, est barrée, et la page s'arrête
sur « Mais impliquent-elles (iv) ? ».

\begin{quote}
\textbf{Définition} (page 85). Quand ces conditions sont satisfaites, $F'$ est une
\emph{subdivision} de $F$, et $F$ une \emph{surdivision} de $F'$ ; on écrit $F'
\preccurlyeq F$.
\end{quote}

C'est un ordre : la transitivité suit de (i). Comme $G \parallel F$ entraîne
toujours $G \parallel F'$ par (1.58), la condition (i) dit que $F'$ et $F$ ont
les mêmes figures disjointes, et l'on retrouve (1.25) de la page 12. Une
subdivision de $F$ induit une subdivision de toute sous-figure, par le lemme de
la page 82 ; et réciproquement, si $F = \mathrm{Sup}\,F_{i}$, $F'$ est une
subdivision de $F$ dès que chaque $F'|F_{i}$ en est une de $F_{i}$, puisque
l'ombre de $F$ est la réunion de celles des $F_{i}$\footnote{La page 86 conclut
« ce qui par le lemme signifie $x \parallel F'$ », sans la barre de négation
qu'appelle le raisonnement : $x \nparallel F'$.}.

\begin{quote}
\textbf{At~14} (page 86, \emph{recollement des subdivisions}). \emph{Si $G \leq F$
et $G' \preccurlyeq G$, il existe une subdivision de $F$ qui induit $G'$.}
\end{quote}

Par recollement --- transfini si $\widetilde{F}$ est infini --- on doit se ramener
au cas où $F$ est une multistrate. On aimerait de plus un choix canonique, la
moins fine des subdivisions voulues, « comme dans le cas ensembliste » :

\begin{quote}
\textbf{At (fort) 14 bis} (page 87). \emph{On peut choisir $F' \preccurlyeq F$
induisant $G'$ de sorte que, $\varphi : \widetilde{F'} \to \widetilde{F}$ étant
l'application canonique, le carré}
\end{quote}

\begin{tikzcd}
\widetilde{G'} \arrow[r, hook] \arrow[d, "\varphi|\widetilde{G'}"'] & \widetilde{F'} \arrow[d, "\varphi"] \\
\widetilde{G} \arrow[r, hook] & \widetilde{F}
\end{tikzcd}

\begin{quote}
\emph{soit cartésien, que $\varphi$ induise un isomorphisme d'ensembles ordonnés
du complémentaire de $\widetilde{G'_{+}}$ sur $\widetilde{F} \smallsetminus
\widetilde{G}$, et que l'ordre de $\widetilde{F'}$ entre ces deux morceaux soit
l'image réciproque de celui de $\widetilde{F}$.}
\end{quote}

Le 22 juin, il pense qu'un tel $F'$ doit être unique, comme la plus grande des
subdivisions de $F$ induisant $G'$, et que, pour des figures finies, la condition
équivaut à :

\begin{quote}
\textbf{At (fort) 14 ter.} \emph{Pour $X \in \mathcal{M}$ et $F' \preccurlyeq
\partial X$, il existe une unique $X' \in \mathcal{M}$ avec $X' \preccurlyeq X$ et
$\partial X' = F'$.}
\end{quote}

C'est ce que font les cellules : une boule dont on subdivise le bord reste une
seule cellule. Ce n'est pas vrai dans un atelier simplicial strict, où chaque
multistrate a la combinatoire d'un simplexe et l'intersection de deux
multistrates compatibles est une multistrate : un triangle dont un côté a reçu un
sommet supplémentaire n'est plus le bord d'aucun simplexe\footnote{La page 88
dessine le triangle, un point marqué près d'un sommet ; la phrase qui
l'accompagne est barrée et lue en partie, et la lecture du contre-exemple est la
nôtre. L'indice $+$ de $\widetilde{G'_{+}}$ et l'étiquette de la flèche de gauche
sont douteux. Le rapprochement avec les complexes cellulaires réguliers est
nôtre.}. Il imposera du moins At~14 ter aux ateliers \emph{stricts}, à définir
plus bas.

\pagerange{88}{93}
\section*{Le cas ensembliste, et les sous-ateliers (pages 88 à 93)}

La page 88 fixe un programme : expliciter les axiomes pour un atelier
ensembliste ; étudier leur stabilité par passage à un sous-atelier ; étudier la
\emph{saturation} d'un atelier. Les deux premiers points occupent les dernières
pages ; le troisième n'est pas abordé.

\emph{Le cas ensembliste.} Un atelier ensembliste, donné par $\mathcal{L}$ et
$\mathfrak{F} \subset \mathrm{Fig}(\mathcal{L})$, satisfait automatiquement At~1 à
At~8. Lorsque $\mathcal{M}$ est formé des seuls points $\lbrace \lbrace x \rbrace
\rbrace$, la donnée se réduit à une relation réflexive et symétrique sur
$\mathcal{L}$. En général, les axiomes se traduisent ainsi :
\begin{enumerate}
\item[At~9] devient, compte tenu de Atens~2 : \emph{deux multistrates de supports
disjoints sont compatibles}, donc disjointes ; deux multistrates sont alors
disjointes si et seulement si leurs supports le sont.
\item[At~9$_{\Sigma}$] devient Atens~2$_{\Sigma}$ : \emph{$X
\between_{\mathcal{M}} Y$ si et seulement si $X$ et $Y$ sont ensemblistement
compatibles, c'est-à-dire $|X| \cap |Y|$ est la réunion $K$ de leurs strates
communes, et $\widetilde{K} \in \Sigma$} ; pour $\Sigma$ formé de tous les
ensembles ordonnés, c'est dire que $\mathfrak{F}$ est un sous-atelier spécieux de
$\mathrm{Fig}(\mathcal{L})$, ce qui n'arrive presque jamais pour les ateliers
simpliciaux.
\item[At~10] est automatique, par Atens~2.
\item[At~11 bis] est vide pour des figures de type fini, la famille étant alors
stationnaire.
\item[At~12] devient Atens~4 : \emph{si $X \between_{\mathcal{M}} Y$, $K = X \cap
Y$, $X' \ll X$, $Y' \ll Y$ et $X'|K \between_{\mathcal{M}} Y'|K$, alors $X'
\between_{\mathcal{M}} Y'$} ; il résulte de Atens~2$_{\Sigma}$\footnote{La page
conclut « alors $X \between Y$ » ; la traduction de At~12 demande $X' \between
Y'$.}.
\item[At~13] signifie que $F \ll X$ avec $|F| \subsetneq |X|$, et il est toujours
vérifié.
\item[At~14 bis] devient Atens fort~5 : \emph{si $X \in \mathcal{M}$ et $F'
\preccurlyeq \partial X$, l'unique figure ensembliste élémentaire de bord $F'$ et
de support $|X|$ est dans $\mathcal{M}$} --- « mais ce sera pas vrai pour les
ateliers simpliciaux ! ». At~14, en revanche, se traduit mal en termes de
sous-ateliers : ce n'est plus une condition de stabilité.
\end{enumerate}

\emph{Sous-ateliers} (pages 91 à 93). Pour $\mathfrak{F}' \subset \mathfrak{F}$
satisfaisant Sousat~1 et Sousat~2\footnote{La page 91 renvoie aux conditions
« Sous at 1, 2, p.~47 », chiffre douteux ; elles sont à la page 42.}, At~1 à
At~7 passent automatiquement. Pour les axiomes où entrent les lieux, il faut
distinguer $\mathcal{L}$ et $\mathcal{L}'$ : on a $\mathcal{L} \cap \mathfrak{F}'
\subset \mathcal{L}'$, et l'on voudrait l'égalité. D'où :
\begin{enumerate}
\item[Sous-at~3] (\emph{épais}) : si $X \in \mathcal{M}' = \mathfrak{F}' \cap
\mathcal{M}$ et $x \in \mathcal{L}$, $x \ll X$, alors $x \in \mathfrak{F}'$ ;
\item[Sous-at~3'] (\emph{quasi-dense}) : pour tout $X \in \mathcal{M}'$, il existe
$x \in \mathcal{L} \cap \mathfrak{F}'$ avec $x \mathrel{\overset{\circ}{\ll}} X$ ;
\item[Sous-at~4] : si $X, Y \in \mathcal{M}'$ sont disjoints dans $\mathfrak{F}$,
ils le sont dans $\mathfrak{F}'$, c'est-à-dire $X \vee Y \in \mathfrak{F}'$.
\end{enumerate}
Sous-at~3 entraîne Sous-at~3' (par At~8 dans $\mathfrak{F}$) ; chacune entraîne
At~8 pour $\mathfrak{F}'$, et la seconde déjà $\mathcal{L}' = \mathcal{L} \cap
\mathfrak{F}'$. At~9 ne passe pas nécessairement à un sous-atelier, même plein ;
il passe moyennant Sous-at~3 et Sous-at~4, ainsi que At~10\footnote{La page 93
écrit que At~10 pour $\mathfrak{F}'$ implique At~10 pour $\mathfrak{F}$ ; le sens
qu'appelle le reste de la page est de $\mathfrak{F}$ vers $\mathfrak{F}'$.} et
At~13. At~11 bis et At~12 sont des propriétés de stabilité de $\mathfrak{F}'$ dans
$\mathfrak{F}$. At~14 reste « discordante » --- lecture douteuse --- et At~14 bis
est, elle, une condition de stabilité. Un trait clôt la page 93, le chapitre et
le dossier.

\end{document}
